% Standalone expansion; regenerate with python3 code/build_article.py
\documentclass[11pt,letterpaper]{article}
\usepackage[T1]{fontenc}
\usepackage[utf8]{inputenc}
\usepackage{lmodern}
\usepackage[sc]{mathpazo}
\usepackage{amsmath,amssymb,amsthm,mathtools,bm}
\usepackage[margin=1in,headheight=15pt]{geometry}
\usepackage{microtype,graphicx,xcolor,booktabs,longtable,tabularx,array}
\usepackage{enumitem,listings,fancyhdr}
\usepackage[hyphens]{url}
\usepackage{hyperref,bookmark}
\definecolor{Ink}{HTML}{17344A}
\definecolor{Accent}{HTML}{24655A}
\definecolor{Muted}{HTML}{52616A}
\hypersetup{colorlinks=true,linkcolor=Ink,citecolor=Accent,urlcolor=Accent,
 pdftitle={Twisted Stieltjes Convolutions and Harmonic Laurent Identities},
 pdfauthor={Research continuation prepared for Vladimir Reshetnikov},
 pdfsubject={Polylogarithms, Gamma, polygamma, generalized harmonic numbers, Hurwitz zeta and Stieltjes constants}}
\pagestyle{fancy}
\fancyhf{}
\fancyhead[L]{\small\textcolor{Muted}{Twisted convolutions and harmonic Laurent identities}}
\fancyhead[R]{\small\textcolor{Muted}{ProveIt}}
\fancyfoot[C]{\thepage}
\renewcommand{\headrulewidth}{0.3pt}
\setlength{\parindent}{1.2em}
\setlength{\parskip}{3pt}
\setlength{\emergencystretch}{2em}
\allowdisplaybreaks[2]
\numberwithin{equation}{section}
\theoremstyle{plain}
\newtheorem{theorem}{Theorem}[section]
\newtheorem{proposition}[theorem]{Proposition}
\newtheorem{lemma}[theorem]{Lemma}
\newtheorem{corollary}[theorem]{Corollary}
\theoremstyle{definition}
\newtheorem{definition}[theorem]{Definition}
\newtheorem{question}[theorem]{Question}
\newtheorem{conjecture}[theorem]{Conjecture}
\theoremstyle{remark}
\newtheorem{remark}[theorem]{Remark}
\newcommand{\bbT}{\mathbb T}
\newcommand{\bbC}{\mathbb C}
\newcommand{\bbR}{\mathbb R}
\newcommand{\bbZ}{\mathbb Z}
\newcommand{\bbQ}{\mathbb Q}
\newcommand{\fp}{\operatorname{FP}}
\newcommand{\res}{\operatorname{Res}}
\newcommand{\numi}{\mathrm i}
\newcommand{\rd}{\mathrm d}
\newcommand{\rme}{\mathrm e}
\newcommand{\rLi}{\operatorname{Li}}
\newcommand{\ordder}[2]{\left.\partial_s^{#1}\rLi_s(#2)\right|_{s=0}}
\newcommand{\pinned}{\texttt{e0d9463bdee9685dfb1dddb819059cc738540c57}}
\begin{document}
\begin{titlepage}
\thispagestyle{empty}
{\color{Accent}\large PROVEIT RESEARCH CONTINUATION\par}
\vspace{1.25cm}
{\color{Ink}\huge\bfseries Twisted Stieltjes Convolutions\\[5pt]
and Harmonic Laurent Identities\par}
\vspace{0.65cm}
{\Large Periodic Stieltjes constants, polygamma contact terms,\\
polylogarithm order derivatives, and shifted harmonic sums\par}
\vspace{0.8cm}
{\large Prepared for Vladimir Reshetnikov\par}
{\small Research continuation prepared with OpenAI assistance\par}
{\large 10 October 2026\par}
\vspace{0.7cm}
\noindent\textcolor{Accent}{\rule{\textwidth}{0.6pt}}
\begin{abstract}
We continue the exact-identity programme of the ProveIt polylogarithm
manuscript in two directions. First, we answer an archived research
question about the all-index half-twist Stieltjes multiplication law,
including point-supported counterterms. A meromorphic periodic Lerch
family gives the construction for every real nonintegral twist. Its
convolution identity is the full Dirac mass, and its covariant
derivative correction is an explicit harmonic-number multiple of a
Dirac derivative. Rational twists reduce to finite Hurwitz and
Stieltjes combinations. At the half-twist, an ordinary convergent
subtracted digamma integral evaluates in terms of
\(\log\Gamma(1/4)\), Euler's constant, and logarithms.
An exact zero-frequency subtraction controls the untwisted limit.

Second, a two-variable Mellin continuation determines every principal
Laurent coefficient and finite part of
\(\sum_{n\geq1}H_{n-1}(\{1\}^r)/(n+a)^s\)
at \(s=1\) and at all nonpositive integers. Gamma and generalized
Bernoulli coefficients give a finite formula at every harmonic depth
and complex outer shift. We derive residue-aware differentiation,
depth-mixing reflection, rational multiplication, and complete
antiderivative ladders for these coefficient polynomials, with
explicit half-shift examples.

For self-contained proofs and consistent normalizations we consolidate
the untwisted periodic Stieltjes algebra and polylogarithm kernels
already developed in archived repository continuations, crediting those
results and their classical antecedents. The package includes exact
symbolic checks, independent numerical diagnostics, a source audit,
two proposed manuscript corrections, and further research questions.
No universal priority claim is made; the remaining \(S_6\) and newer
\(S_8\) special-value conjectures remain unresolved.
\end{abstract}
\vfill
\noindent\small
Source baseline: \pinned.\par
The archive contains this article, its complete \LaTeX\ source,
reproducible verification code, a source manifest, and proposed
integration patches.
\end{titlepage}
\tableofcontents
\clearpage

% BEGIN sections/01_scope.tex
\section{Scope, conventions, and the status of the results}
\label{scope:sec}

This continuation develops two exact extensions of the ProveIt programme.
First, it answers the archived question about an all-index half-twist
Stieltjes multiplication law with its endpoint terms. The twist changes
the convolution identity from \(\delta_0-1\) to \(\delta_0\).
Second, it determines the principal Laurent coefficients and finite
parts of a shifted harmonic Dirichlet family at every nonpositive
integer, with reflection, rational multiplication, and antiderivative
formulas for those coefficients.

Both developments require a precise normalization at a moving pole.
The periodic correction is supported at the endpoint of a circle.
The harmonic correction is a Laurent residue, which survives when
finite parts are differentiated. Their common role explains why a
formula at regular parameters cannot simply be specialized or
differentiated at a singular parameter.

\subsection{The inspected baseline}

We froze the repository at \pinned, whose commit is dated
10 October 2026, 21:44:12 UTC. The audit used the complete set of
72 chapter sources and the eight non-PDF manuscript support files.
The pending archive
\path{ProveIt_Polylogarithm_Stieltjes_Continuation_2026-10-10.zip}
was matched to its Git blob and its mathematical source was inspected.
Five relevant archived report sources were also checked.
The machine-readable source manifest records exact paths and hashes.
These statements describe the inspected set; they do not assert that
every report or historical archive in the repository was reviewed.
Manuscript and archival status are kept separate
\cite{RepoManuscript,RepoIncoming}.

The archived \emph{Convolution Algebra} report explicitly asks for a
nonintegral-twist normal form and, as a first target, the all-index
half-twist multiplication table with point-supported counterterms
\cite[final research section]{RepoConvolutionAlgebra}.
Section~\ref{tw:sec:main} supplies this normalization, multiplication
law, and covariant derivative/primitive calculus. It gives a
rational-twist conversion to Hurwitz and Stieltjes coordinates, a
convergent half-twist integral, and the exact zero-frequency correction
needed in the limit as the twist tends to zero.

The same archive and two further continuations already contain the
untwisted periodic Bell algebra, polygamma contact terms, normalized
primitive tower, and reflected polylogarithm kernels
\cite{RepoConvolutionAlgebra,RepoConvolutionCalculus,RepoShiftedJets}.
Those results are consolidated with independent complete proofs here;
they are not new theorems relative to the pinned repository.
The incoming Stieltjes continuation supplies additional ordinary
translation and integration formulas \cite{RepoIncoming}.
Earlier alternating-harmonic reports likewise already use beta
representations, cyclotomic projections, and scalar Mellin subtraction
\cite{RepoAlternating,RepoHarmonicGeometry}.

The source's \(S_2\) and \(S_4\) identities already have exact proofs.
The frozen \(S_6\) target and the newer \(S_8\) candidate remain
conjectural. Our proofs do not evaluate either of them.
The half-twist target answered here is an explicit research question;
it is distinct from those two special-value conjectures.

\subsection{Classical inputs and the contribution of this article}

The Fourier coefficients of Hurwitz zeta, Gamma--beta identities, and
\[
 \zeta'(0,x)=\log\Gamma(x)-\frac12\log(2\pi)
\]
are classical \cite{DLMFHurwitz,DLMFBeta}.
Normalized Hurwitz convolution is established; Sun states an
ordinary-integral form in his study of invariant functions
\cite[Remark 3.2 in the arXiv version]{per:Sun}.
Periodic complex-power distributions provide the natural background
\cite{per:GaySebbar}. Kim's Hurwitz-transform continuation is a recent
antecedent for scalar test-function subtraction \cite{Kim2026}.

Young's parametric Euler sums supply the beta-derivative evaluation
framework \cite{Young2026}. His height-one work and unshifted Laurent
study are directly relevant \cite{Young2024,Young2023}; the harmonic
Stieltjes literature concerns related, but not identically shifted,
families \cite{HarmonicStieltjes2023}. The full text of Young's 2023
paper was unavailable. Its inspected publisher abstract already rules
out a broad priority claim for unshifted Laurent expansions.

The additions for this continuation are the fully normalized twisted
extension and the joint coefficient theorem for the outer-shifted
harmonic family, with its finite polynomial calculus.
Beta differentiation, Mellin continuation, and the unshifted height-one
theory are treated as antecedents. Some consequences may have
equivalent formulations elsewhere. A result described as \emph{proved
here} has a complete analytic proof in this article; a result described
as an \emph{addition to the inspected source} was not located in that
source set. Neither phrase asserts universal literature priority.
Proof and provenance statuses are recorded separately.

\begin{center}
\begin{tabularx}{\textwidth}{@{}p{0.25\textwidth}X@{}}
\toprule
Result or ingredient & Status and use\\
\midrule
Untwisted periodic calculus &
Existing repository foundation, independently checked; all endpoint
terms and integration constants retained.\\[3pt]
Twisted Stieltjes calculus &
Answers the archived half-twist target, for every real nonintegral
twist, with the singular normalization included.\\[3pt]
Ordinary half-twist integral &
A convergent consequence with a midpoint evaluation in
\(\log\Gamma(1/4)\), Euler's constant, and logarithms.\\[3pt]
Harmonic beta evaluations &
Established antecedent connecting harmonic polynomials, Gamma
derivatives, and positive integer values.\\[3pt]
Shifted harmonic Laurent data &
Complete principal and finite parts; residue-aware differentiation,
reflection, rational multiplication, and repeated primitives.\\[3pt]
Numerical checks &
Independent diagnostics; none replaces an analytic proof.\\
\bottomrule
\end{tabularx}
\end{center}

\subsection{Notation and normalizations}

Order derivatives always concern the first argument:
\[
 \zeta^{[j]}(s,a)=\partial_s^j\zeta(s,a),\qquad
 \rLi_s^{[j]}(z)=\partial_s^j\rLi_s(z).
\]
By contrast, \(\psi^{(k)}(a)\) differentiates the argument \(a\), and
\begin{equation}\label{scope:stieltjes}
 \zeta(1+w,a)=\frac1w+
       \sum_{n\geq0}\frac{(-1)^n}{n!}\gamma_n(a)w^n.
\end{equation}
Thus \(\gamma_0(a)=-\psi(a)\), and \(\gamma=\gamma_0(1)\).
Finite parts are Laurent constant terms, unless a displayed cutoff
definition specifies the equivalent integral normalization.

On \(\bbT=\bbR/\bbZ\), write
\[
 \widehat T(n)=\langle T,\rme^{-2\pi\numi n x}\rangle,\qquad
 \Delta=\delta_0-1.
\]
The distribution \(\Delta\) is the identity for convolution on the
subspace of distributions with zero mean. The full algebra has identity
\(\delta_0\). Reflection means
\(\check T(x)=T(-x)\), and \(D\) denotes the distributional argument
derivative. A convolution is always the compact-circle convolution.
It is defined for any pair of distributions on the circle; it must
not be confused with a pointwise product of singular distributions.
All logarithms of \(2\pi\numi n\), for \(n\ne0\), mean
\[
 \log(2\pi\numi n)=\log(2\pi|n|)
                 +\frac{\pi\numi}{2}\operatorname{sgn}(n).
\]
This branch convention is part of the identities.

For finite harmonic sums we use
\[
 H_N^{(j)}=\sum_{k=1}^N k^{-j},\qquad
 H_N(\{1\}^r)=
 \sum_{N\geq n_1>\cdots>n_r\geq1}\frac1{n_1\cdots n_r},
 \qquad H_N(\{1\}^0)=1.
\]
Their generating polynomial is
\[
 \sum_{r\geq0}H_N(\{1\}^r)u^r
      =\prod_{k=1}^N(1+u/k).
\]
Thus all depths concern explicit polynomials in generalized harmonic
numbers:
\[
 H_N(\{1\}^2)=\frac{H_N^2-H_N^{(2)}}2,\qquad
 H_N(\{1\}^3)=\frac{H_N^3-3H_NH_N^{(2)}+2H_N^{(3)}}6.
\]

% END input


% BEGIN sections/02_periodic_calculus.tex
\section{Periodic finite parts and harmonic contact terms}
\label{per:sec:main}

We begin by consolidating the untwisted foundation already present in
three archived ProveIt continuations
\cite{RepoConvolutionAlgebra,RepoConvolutionCalculus,RepoShiftedJets}.
The endpoint normalization is fixed by analytic continuation from the
integrable Hurwitz family. Its spectral finite parts and argument
derivatives differ by an explicit harmonic-number multiple of a Dirac
mass derivative. This distinction will persist in the twisted extension.

The normalized Hurwitz convolution is classical
\cite[Remark 3.2]{per:Sun}; complex-power distributions provide related
background \cite[Section 5]{per:GaySebbar}. Scalar meromorphic
continuation after pairing with a smooth function is also covered by
Kim's Hurwitz-transform theorem
\cite[Theorem 1 and Corollary 2]{Kim2026}. The full proofs below establish
the conventions used in the rest of this article and independently check
the archived formulas; the untwisted Bell algebra and primitive tower are
not new-to-repository claims.

\subsection{Conventions and the meromorphic periodic family}

Write $\mathbb T=\mathbb R/\mathbb Z$, use normalized Lebesgue measure,
and put
\[
 \widehat U(n)=\langle U,e^{-2\pi i n x}\rangle,\qquad
 D=\frac{d}{dx},\qquad \Delta=\delta_0-1.
\]
The Dirac mass has Fourier coefficient $1$ at every integer. Thus
$\widehat\Delta(0)=0$ and $\widehat\Delta(n)=1$ for $n\ne0$.
On the convolution algebra of distributions with zero mean, $\Delta$
is the identity. Reflection is $\check U(x)=U(-x)$, so
$\widehat{\check U}(n)=\widehat U(-n)$. Throughout,
\begin{equation}\label{per:eq:L}
 L_n=\log(2\pi |n|)+\frac{\pi i}{2}\operatorname{sgn}n,
 \qquad n\ne0.
\end{equation}
Consequently $(2\pi i n)^z$ always means $\exp(zL_n)$.

\begin{theorem}[The entire normalized family and all residues]
\label{per:thm:family}
There is an entire distribution-valued function $W_s$ characterized by
\begin{equation}\label{per:eq:W}
 \widehat W_s(0)=0,\qquad
 \widehat W_s(n)=(2\pi i n)^{s-1}\quad(n\ne0).
\end{equation}
The family
\begin{equation}\label{per:eq:Z}
 Z_s=\Gamma(1-s)W_s
\end{equation}
is the unique meromorphic continuation, as a periodic distribution, of
$x\mapsto\zeta(s,x)$ from $\operatorname{Re}s<1$. Its poles are simple
and occur at every positive integer. Their residues are
\begin{align}
 \operatorname*{Res}_{s=1} Z_s&=1-\delta_0=-\Delta,
 \label{per:eq:res1}\\
 \operatorname*{Res}_{s=k+1} Z_s
 &=\frac{(-1)^{k+1}}{k!}\delta_0^{(k)}\qquad(k\ge1).
 \label{per:eq:resk}
\end{align}
For all complex $s,t$, the identities
\begin{equation}\label{per:eq:group}
 DW_s=W_{s+1},\qquad W_s*W_t=W_{s+t-1}
\end{equation}
hold. Equivalently, as meromorphic distribution identities,
\begin{equation}\label{per:eq:meroconv}
 DZ_s=-sZ_{s+1},\qquad
 Z_s*Z_t=\frac{\Gamma(1-s)\Gamma(1-t)}{\Gamma(2-s-t)}Z_{s+t-1}.
\end{equation}
At intersections of polar divisors the second expression is interpreted
by its equal Gamma-normalized expression, before coefficients are taken.
\end{theorem}

\begin{proof}
On every compact set of $s$ the coefficients in \eqref{per:eq:W} have
uniform polynomial growth in $|n|$. Every fixed order derivative in $s$
adds only a power of $L_n$. Since the Fourier coefficients of a smooth
test function decrease faster than every inverse power, the series
\[
 \langle W_s,\phi\rangle
 =\sum_{n\ne0}(2\pi i n)^{s-1}\widehat\phi(-n)
\]
and all its fixed order derivatives converge normally on compact sets.
This constructs the entire family. Hurwitz's classical Fourier formula
\cite[\S25.11]{DLMFHurwitz} identifies \eqref{per:eq:Z} with $\zeta(s,x)$
for $\operatorname{Re}s<0$. The relation
$\zeta(s,x)=x^{-s}+\zeta(s,1+x)$ shows that the latter is holomorphic
with values in $L^1(0,1)$ on $\operatorname{Re}s<1$, so the
identification extends throughout that half-plane. The identity theorem
for each test-function pairing gives uniqueness of the continuation.

At $s=k+1+w$,
\begin{equation}\label{per:eq:gammapole}
 \Gamma(-k-w)=\frac{(-1)^{k+1}}{k!w}
       +\frac{(-1)^k}{k!}(H_k-\gamma)+O(w),
 \qquad H_0=0.
\end{equation}
The Fourier coefficients show that $W_1=\Delta$ and
$W_{k+1}=D^k\Delta=\delta_0^{(k)}$ for $k\ge1$.
This proves the residues and also shows that none of the listed poles
is removable. Multiplication of Fourier coefficients proves
\eqref{per:eq:group}. Convolution of two distributions on the compact
group $\mathbb T$ is well defined. Normal convergence after pairing with
a test function justifies the parameter identities jointly. Finally,
$\Gamma(1-s)=-s\Gamma(-s)$ proves the derivative identity in
\eqref{per:eq:meroconv}.
\end{proof}

The residue at $s=1$ has both a constant and a Dirac term. The constant
is the ordinary Hurwitz pole seen on every open subinterval of $(0,1)$;
the Dirac term comes from the nonintegrable endpoint $x^{-s}$. Omitting
either term destroys the zero Fourier coefficient. The poles at
$2,3,\ldots$ are entirely supported at the endpoint and are invisible
in the pointwise meromorphic function of $s$ for a fixed $0<x<1$.

\subsection{Cutoff definitions and the normalization at the endpoint}

Define $T_n$ by the Laurent expansion
\begin{equation}\label{per:eq:Tdef}
 Z_{1+w}=-\frac{\Delta}{w}
       +\sum_{n=0}^{\infty}\frac{(-1)^n}{n!}T_nw^n.
\end{equation}
Each $T_n$ agrees with $\gamma_n(x)$ on $0<x<1$. Define also
\begin{equation}\label{per:eq:Pdef}
 F_k=\operatorname{FP}_{s=k+1}Z_s,\qquad
 P_k=(-1)^{k+1}k!F_k\quad(k\ge0),
\end{equation}
where $\operatorname{FP}$ denotes the constant Laurent coefficient.
Thus $F_0=T_0$ and $P_0=-T_0$. For $k\ge1$, the usual
$\psi^{(k)}(x)=(-1)^{k+1}k!\zeta(k+1,x)$ shows that $P_k$ agrees with
the ordinary $k$th polygamma function away from $0$.

\begin{proposition}[Equivalent Hadamard finite parts]
\label{per:prop:cutoffs}
For every smooth periodic test function $\phi$,
\begin{equation}\label{per:eq:Tcutoff}
 \langle T_n,\phi\rangle
 =\lim_{\epsilon\downarrow0}\left\{
   \int_\epsilon^1\gamma_n(x)\phi(x)\,dx
    +\frac{\phi(0)}{n+1}\log^{n+1}\epsilon\right\}.
\end{equation}
For $k\ge1$,
\begin{align}
 \langle F_k,\phi\rangle
 =\lim_{\epsilon\downarrow0}\biggl\{&
   \int_\epsilon^1\zeta(k+1,x)\phi(x)\,dx
   -\sum_{j=0}^{k-1}\frac{\phi^{(j)}(0)}{j!}
                       \frac{\epsilon^{j-k}}{k-j}
   +\frac{\phi^{(k)}(0)}{k!}\log\epsilon\biggr\}.
 \label{per:eq:Fcutoff}
\end{align}
All these distributions have zero mean.
\end{proposition}

\begin{proof}
Subtract a Taylor polynomial of the test function at $0$. For any
$N\ge0$, the pairing of $Z_s$ has the continuation
\begin{align}
 \langle Z_s,\phi\rangle={}&
 \int_0^1x^{-s}\left(\phi(x)-\sum_{j=0}^{N}
                         \frac{\phi^{(j)}(0)}{j!}x^j\right)dx
 +\sum_{j=0}^{N}\frac{\phi^{(j)}(0)}{j!(j+1-s)}\notag\\
 &+\int_0^1\zeta(s,1+x)\phi(x)\,dx,
 \label{per:eq:localcontinuation}
\end{align}
in $\operatorname{Re}s<N+2$, apart from the displayed and Hurwitz
poles. Taking the constant coefficient at $s=k+1$, with $N=k$,
proves \eqref{per:eq:Fcutoff} by elementary integration of the Taylor
monomials. At $s=1+w$, use
\[
 \gamma_n(x)=\frac{\log^n x}{x}+\gamma_n(1+x).
\]
The coefficient in \eqref{per:eq:localcontinuation} is
\[
 \int_0^1\frac{\log^n x}{x}\bigl(\phi(x)-\phi(0)\bigr)dx
 +\int_0^1\gamma_n(1+x)\phi(x)\,dx.
\]
This is precisely \eqref{per:eq:Tcutoff}. The zero-mean assertion
follows from \eqref{per:eq:W} by taking Laurent coefficients.
\end{proof}

Zero mean alone does not specify a periodic extension of a singular
function: derivatives of a Dirac mass have zero mean as well. Here the
normalization is fixed by \eqref{per:eq:Z}, or equivalently by the
cutoff formulas. This qualification is essential in using the word
``canonical.''

\subsection{The harmonic contact correction}

\begin{theorem}[Spectral finite parts and compatible derivatives]
\label{per:thm:contact}
For $n\ne0$ and $k\ge0$,
\begin{equation}\label{per:eq:Pfourier}
 \widehat P_k(n)=(2\pi i n)^k\bigl(\gamma+L_n-H_k\bigr),
 \qquad \widehat P_k(0)=0.
\end{equation}
Consequently,
\begin{equation}\label{per:eq:anomaly}
 DP_k=P_{k+1}+\frac{1}{k+1}\delta_0^{(k+1)}.
\end{equation}
If
\[
 g(x)=\log\Gamma(x)-\frac12\log(2\pi),\quad 0<x<1,
\]
is periodized as a locally integrable function, then $Dg=P_0$. The
derivative-compatible extensions
\begin{equation}\label{per:eq:Q}
 Q_k=D^{k+1}g=D^kP_0
\end{equation}
satisfy $DQ_k=Q_{k+1}$ exactly, and
\begin{equation}\label{per:eq:Qcorrection}
 Q_k=P_k+H_k\delta_0^{(k)}\qquad(k\ge1),\qquad Q_0=P_0.
\end{equation}
\end{theorem}

\begin{proof}
Expand $W_{k+1+w}=W_{k+1}+w\partial_sW_s|_{s=k+1}+O(w^2)$
against \eqref{per:eq:gammapole}; after multiplication by
$(-1)^{k+1}k!$, the finite Fourier coefficient is
$(2\pi i n)^k(\gamma+L_n-H_k)$.
Since $H_{k+1}-H_k=1/(k+1)$, differentiation of this coefficient gives
\eqref{per:eq:anomaly}. Lerch's identity gives $g=\partial_sZ_s|_{s=0}$.
Taking the first coefficient in $DZ_s=-sZ_{s+1}$ at $s=0$ gives
$Dg=-T_0=P_0$; the coefficient at degree zero separately says
$DZ_0=\Delta$. Finally the Fourier coefficient of $D^kP_0-P_k$ is
$H_k(2\pi i n)^k$, which is \eqref{per:eq:Qcorrection}.
\end{proof}

For example, the extension given by the finite part of
$\zeta(s,x)$ at $s=2$ yields a periodic trigamma $P_1$, whereas
$D^2g=P_1+\delta_0'$. Both agree with $\psi'(x)$ on $(0,1)$.
The discrepancy is supported precisely at the identified endpoints;
it is not a failure of ordinary differentiation on the open interval.


% END input


% BEGIN sections/03_convolutions.tex
\section{Stieltjes convolution identities and repeated primitives}
\label{per:sec:convolutions}

\subsection{All Stieltjes functions as convolution polynomials}

The regularized germ
\begin{equation}\label{per:eq:B}
 B(u)=-uZ_{1+u}=\Gamma(1-u)W_{1+u},\qquad |u|<1,
\end{equation}
satisfies
\begin{equation}\label{per:eq:Bseries}
 B(u)=\Delta+
       \sum_{r\ge1}\frac{(-1)^r}{(r-1)!}T_{r-1}u^r.
\end{equation}
The following formula is a statement in the commutative convolution
algebra with identity $\Delta$; exponential and Bell-polynomial
expressions can equivalently be read as formal power series.

\begin{theorem}[Stieltjes convolution generator]
\label{per:thm:Bell}
For $|u|<1$,
\begin{equation}\label{per:eq:exp}
 B(u)=\exp_*\left(uP_0+
                    \sum_{r\ge2}\frac{\zeta(r)}{r}u^r\Delta\right).
\end{equation}
Here the series defines the analytic germ after pairing with any smooth
test function. In particular, if $\mathcal B_j^*$ denotes the complete
exponential Bell polynomial with every multiplication interpreted as
convolution and every scalar constant as that scalar times $\Delta$,
then
\begin{equation}\label{per:eq:Bell}
 T_n=\frac{(-1)^{n+1}}{n+1}
 \mathcal B_{n+1}^*
 \bigl(P_0,\zeta(2)\Delta,2!\zeta(3)\Delta,
                         \ldots,n!\zeta(n+1)\Delta\bigr).
\end{equation}
The two-parameter form is
\begin{equation}\label{per:eq:Bconv}
 B(u)*B(v)=C(u,v)B(u+v),
 \qquad C(u,v)=\frac{\Gamma(1-u)\Gamma(1-v)}{\Gamma(1-u-v)},
\end{equation}
as an analytic germ at $(u,v)=(0,0)$. Its scalar correction has
\begin{equation}\label{per:eq:C}
 \log C(u,v)=\sum_{r\ge2}\frac{\zeta(r)}{r}
                   \bigl(u^r+v^r-(u+v)^r\bigr).
\end{equation}
Thus every $T_m*T_n$ is a finite explicit linear combination of
$\Delta,T_0,\ldots,T_{m+n+1}$, with coefficients polynomial in positive
integer zeta values.
\end{theorem}

\begin{proof}
At a nonzero Fourier mode, $\widehat P_0(n)=\gamma+L_n$. The Taylor
expansion
\[
 \log\Gamma(1-u)=\gamma u+
                     \sum_{r\ge2}\frac{\zeta(r)}r u^r
\]
therefore proves \eqref{per:eq:exp} mode by mode. The same polynomial
growth argument as in Theorem~\ref{per:thm:family} proves normal
convergence in every smaller disk after test-function pairing. Comparing
the complete Bell-polynomial generating function with
\eqref{per:eq:Bseries} gives \eqref{per:eq:Bell}; the factorial is
$n!/(n+1)!=1/(n+1)$. Equation \eqref{per:eq:Bconv} follows from
\eqref{per:eq:group}. The logarithm of its Gamma quotient is
\eqref{per:eq:C}; its linear Euler-constant terms cancel. Coefficient
extraction at $u^{m+1}v^{n+1}$ gives the last assertion.
\end{proof}

The first identities are
\begin{align}
 T_0&=-P_0,\notag\\
 T_1&=\frac12\bigl(P_0^{*2}+\zeta(2)\Delta\bigr),
 \label{per:eq:T1}\\
 T_2&=-\frac13\bigl(P_0^{*3}+3\zeta(2)P_0+2\zeta(3)\Delta\bigr),
 \label{per:eq:T2}\\
 T_3&=\frac14\bigl(P_0^{*4}+6\zeta(2)P_0^{*2}
          +8\zeta(3)P_0+(3\zeta(2)^2+6\zeta(4))\Delta\bigr).
 \label{per:eq:T3}
\end{align}
Here $P_0^{*r}$ is an $r$-fold convolution, not a pointwise power.
Eliminating these powers gives, for example,
\begin{align}
 T_0*T_0&=2T_1-\zeta(2)\Delta,\label{per:eq:00}\\
 T_0*T_1&=\frac32T_2-\zeta(2)T_0+\zeta(3)\Delta,
 \label{per:eq:01}\\
 T_0*T_2&=\frac43T_3-2\zeta(2)T_1+2\zeta(3)T_0
                                      -2\zeta(4)\Delta,
 \label{per:eq:02}\\
 T_1*T_1&=T_3-2\zeta(2)T_1+2\zeta(3)T_0
                                      -\frac14\zeta(4)\Delta.
 \label{per:eq:11}
\end{align}
The last coefficient uses $\zeta(2)^2=\tfrac52\zeta(4)$.

\subsection{The reflected Stieltjes generator}

Define the periodic Hilbert transform on zero-mean distributions by
\begin{equation}\label{per:eq:Hilbert}
 \widehat{\mathcal H U}(n)=-i\operatorname{sgn}(n)\widehat U(n)
 \quad(n\ne0),\qquad \widehat{\mathcal H U}(0)=0.
\end{equation}
Its multiplier is bounded, so it is a continuous operator on periodic
distributions and commutes with the order coefficient operations above.
Reflection in the following theorem acts on the periodic variable only;
it does not conjugate the complex parameters.

\begin{theorem}[All-order reflected convolution]
\label{per:thm:reflectedgenerator}
As an analytic germ at $(u,v)=(0,0)$,
\begin{equation}\label{per:eq:Brefconv}
 B(u)*\check{B(v)}
 =C(u,v)\bigl(\cos(\pi v)\,\mathrm{Id}
                    +\sin(\pi v)\,\mathcal H\bigr)B(u+v),
\end{equation}
with the same Gamma quotient $C(u,v)$ as in \eqref{per:eq:Bconv}.
In particular,
\begin{align}
 T_0*\check T_1={}&\tfrac12T_2+\check T_2
                +\zeta(2)(T_0+\check T_0)+\zeta(3)\Delta,
 \label{per:eq:T0refT1}\\
 T_1*\check T_1={}&\tfrac12(T_3+\check T_3)
                +2\zeta(2)(T_1+\check T_1)
                +\zeta(3)(T_0+\check T_0)
                +\tfrac72\zeta(4)\Delta.
 \label{per:eq:T1refT1}
\end{align}
More generally, every $T_m*\check T_n$ reduces to a finite linear
combination of $\Delta$, $T_j$ and $\check T_j$ with
$0\le j\le m+n+1$ and coefficients polynomial in positive integer
zeta values over $\mathbb Q$.
\end{theorem}

\begin{proof}
Since $L_{-n}=L_n-i\pi\operatorname{sgn}n$, the nonzero Fourier
coefficient of the left side of \eqref{per:eq:Brefconv} is
\[
 \Gamma(1-u)\Gamma(1-v)
       e^{(u+v)L_n}e^{-i\pi v\operatorname{sgn}n}.
\]
The last factor equals
$\cos(\pi v)-i\operatorname{sgn}(n)\sin(\pi v)$, precisely the
multiplier of the operator on the right. The normal-convergence argument
for the normalized family proves the identity as a distribution-valued
analytic germ and justifies every fixed mixed coefficient extraction.
This also checks the sign of the Hilbert-transform term directly.

For explicit coefficients put $b=\gamma+\log(2\pi|n|)$ and
$a=\tfrac{\pi i}{2}\operatorname{sgn}n$, so $a^2=-\pi^2/4$.
Equations \eqref{per:eq:T1}--\eqref{per:eq:T3} give polynomials
$t_j(y)=\widehat T_j(n)$ in $y=b+a$, while reflection replaces $a$
by $-a$. Direct polynomial multiplication gives
\begin{align*}
 t_0(b+a)t_1(b-a)
 ={}&\tfrac12t_2(b+a)+t_2(b-a)
    +\zeta(2)\bigl(t_0(b+a)+t_0(b-a)\bigr)+\zeta(3),\\
 t_1(b+a)t_1(b-a)
 ={}&\tfrac12\bigl(t_3(b+a)+t_3(b-a)\bigr)
    +2\zeta(2)\bigl(t_1(b+a)+t_1(b-a)\bigr)\\
 &+\zeta(3)\bigl(t_0(b+a)+t_0(b-a)\bigr)+\tfrac72\zeta(4).
\end{align*}
Here $\pi^2=6\zeta(2)$ and $\pi^4=90\zeta(4)$ have been used.
These prove the two displayed identities at every nonzero Fourier mode;
all zero modes vanish.

Finally $t_j$ has degree $j+1$ and nonzero rational leading coefficient
$(-1)^{j+1}/(j+1)$. Put $M=m+n+2$. After reducing $a^2=-\pi^2/4$,
the product $t_m(b+a)t_n(b-a)$ has the form
$A(b)+aE(b)$ with $\deg A\le M$ and $\deg E\le M-1$.
The symmetric polynomials $t_j(b+a)+t_j(b-a)$, $0\le j<M$,
have successive degrees $1,\ldots,M$ in $b$, while their antisymmetric
differences divided by $a$ have successive degrees $0,\ldots,M-1$.
Together with the constant they form a triangular basis of this space.
Elimination divides only by nonzero rational leading coefficients.
Its coefficients belong to
$\mathbb Q[\pi^2,\zeta(2),\ldots,\zeta(M)]$, which is the asserted
zeta algebra because $\pi^2=6\zeta(2)$.
\end{proof}

The frequency-sign contribution can equivalently be encoded by
\[
 P_0-\check P_0=-\pi\mathcal H\Delta,\qquad
 (P_0-\check P_0)^{*2}=-\pi^2\Delta.
\]
Thus reflection adds a two-dimensional sign algebra to the unreflected
Bell calculus, with its constant term fixed by the same endpoint mass.

\subsection{Polygamma correlations with all contact terms}

\begin{theorem}[Convolution and reflected convolution]
\label{per:thm:polycorr}
Let $k,l\ge0$ and $r=k+l$. With the finite-part normalization
\eqref{per:eq:Pdef},
\begin{align}
 P_k*P_l=D^r\bigl\{&2T_1+(H_k+H_l)T_0
                         +(H_kH_l-\zeta(2))\Delta\bigr\},
 \label{per:eq:PP}\\
 P_k*\check P_l=(-1)^lD^r\bigl\{&T_1+\check T_1
                   +H_lT_0+H_k\check T_0
                   +(H_kH_l+2\zeta(2))\Delta\bigr\}.
 \label{per:eq:PrefP}
\end{align}
For the derivative-compatible family the harmonic coefficients cancel:
\begin{align}
 Q_k*Q_l&=D^r\{2T_1-\zeta(2)\Delta\},
 \label{per:eq:QQ}\\
 Q_k*\check Q_l&=(-1)^lD^r
              \{T_1+\check T_1+2\zeta(2)\Delta\}.
 \label{per:eq:QrefQ}
\end{align}
\end{theorem}

\begin{proof}
Put $y_n=\gamma+L_n$. The first formula follows from
\[
 \widehat T_0(n)=-y_n,\qquad
 2\widehat T_1(n)=y_n^2+\zeta(2)
\]
and $(y_n-H_k)(y_n-H_l)=y_n^2-(H_k+H_l)y_n+H_kH_l$.
For reflection set $b_n=\gamma+\log(2\pi|n|)$ and
$a_n=\frac\pi2\operatorname{sgn}n$, so $y_n=b_n+ia_n$ and
$y_{-n}=b_n-ia_n$. Then
\[
 (y_n-H_k)(y_{-n}-H_l)
 =b_n^2+a_n^2-(H_k+H_l)b_n+H_kH_l
                                      +ia_n(H_k-H_l).
\]
Also
\[
 \widehat T_1(n)+\widehat T_1(-n)=b_n^2-a_n^2+\zeta(2),
\]
and
$H_l\widehat T_0(n)+H_k\widehat T_0(-n)
=-(H_k+H_l)b_n+ia_n(H_k-H_l)$.
The missing constant is $2a_n^2-\zeta(2)=2\zeta(2)$.
Finally $(2\pi i(-n))^l=(-1)^l(2\pi in)^l$ accounts for the
prefactor in \eqref{per:eq:PrefP}. The last two identities follow
from $Q_k=D^kP_0$ and $\check{D^lP_0}=(-1)^lD^l\check P_0$.
\end{proof}

At the first level,
\begin{equation}\label{per:eq:reflected0}
 P_0*\check P_0=T_1+\check T_1+
                         \frac{\pi^2}{3}(\delta_0-1).
\end{equation}
For $0<a<1$ its ordinary representative is
$\gamma_1(a)+\gamma_1(1-a)-\pi^2/3$. Equation
\eqref{per:eq:reflected0} supplies the missing endpoint mass in the
formal twice-differentiated log-Gamma correlation. In particular, if
$C_g=g*\check g$, then
\begin{equation}\label{per:eq:Cg}
 -D^2C_g=T_1+\check T_1+\frac{\pi^2}{3}(\delta_0-1).
\end{equation}
The negative sign follows from differentiating the reflected factor.

For $r\ge1$, the representative of \eqref{per:eq:PP} on $(0,1)$ can
also be written with one positive-order Hurwitz jet:
\begin{equation}\label{per:eq:PPordinary}
 (P_k*P_l)(a)=(-1)^{r+1}r!
 \left\{2\zeta'(r+1,a)+(2H_r-H_k-H_l)\zeta(r+1,a)\right\}.
\end{equation}
Indeed, differentiating the Hurwitz shift identity at $s=1$ gives
$\gamma_1^{(r)}(a)=(-1)^{r+1}r![\zeta'(r+1,a)+H_r\zeta(r+1,a)]$.
For $r=0$, the separate constant $\zeta(2)$ from $-\zeta(2)\Delta$
must be retained.

Translations cause no further ambiguity. If
$\tau_aU(x)=U(x-a)$ for real $a$ modulo $1$, then
$(\tau_aU)*(\tau_bV)=\tau_{a+b}(U*V)$ and
$\tau_aU*\check{\tau_bV}=\tau_{a-b}(U*\check V)$.
All contact terms above are therefore translated to the sum or difference
of the two shifts. General complex translations are not operators on all
periodic distributions, because their Fourier multipliers grow
exponentially in one direction; they require a separately specified
analytic or one-sided boundary-value setting.

\subsection{Convergent integral identities without distribution notation}

The distribution formulas have ordinary, absolutely convergent integral
versions at every nonintegral shift. The following construction makes
the subtraction completely explicit. For $m,n\ge0$ and $0<a<1$, put
\begin{align}
 \mathcal I_{m,n}(a)={}&
 \int_0^a\left\{\gamma_m(x)\gamma_n(a-x)
       -\gamma_n(a)\frac{\log^m x}{x}
       -\gamma_m(a)\frac{\log^n(a-x)}{a-x}\right\}dx\notag\\
 &+\int_a^1\gamma_m(x)\gamma_n(1+a-x)\,dx\notag\\
 &+\frac{\gamma_n(a)}{m+1}\log^{m+1}a
       +\frac{\gamma_m(a)}{n+1}\log^{n+1}a.
 \label{per:eq:Imn}
\end{align}

\begin{proposition}[Ordinary integral realization]
\label{per:prop:Imn}
Both integrals in \eqref{per:eq:Imn} converge absolutely, and
\[
 \mathcal I_{m,n}(a)=(T_m*T_n)(a),\qquad 0<a<1.
\]
Consequently,
\begin{align}
 \mathcal I_{0,0}(a)&=2\gamma_1(a)+\zeta(2),
 \label{per:eq:I00}\\
 \mathcal I_{0,1}(a)&=\frac32\gamma_2(a)-\zeta(2)\gamma_0(a)-\zeta(3),
 \label{per:eq:I01}\\
 \mathcal I_{0,2}(a)&=\frac43\gamma_3(a)-2\zeta(2)\gamma_1(a)
                         +2\zeta(3)\gamma_0(a)+2\zeta(4),
 \label{per:eq:I02}\\
 \mathcal I_{1,1}(a)&=\gamma_3(a)-2\zeta(2)\gamma_1(a)
                         +2\zeta(3)\gamma_0(a)+\frac14\zeta(4).
 \label{per:eq:I11}
\end{align}
\end{proposition}

\begin{proof}
Near $x=0$, the shift formula for $\gamma_m$ and smoothness of
$\gamma_n$ near $a$ show that the first subtracted integrand is
$O(1+|\log x|^m)$. The same argument applies near $x=a$. The second
integral has no singular endpoints for fixed $a\in(0,1)$.

For a direct identification, first restrict $a$ to a compact subinterval
of $(0,1)$. The singular supports of $T_m(x)$ and $T_n(a-x)$ are
separated in the integration variable. A smooth partition of unity
around $x=0$, $x=a$ and their complement therefore defines their
product: in each singular neighborhood only one factor is a
distribution and the other is smooth. Applying the two cutoff formulas
\eqref{per:eq:Tcutoff}, and then combining the complementary integrals,
gives
\begin{align*}
 \lim_{\epsilon\downarrow0}\biggl\{&
 \int_\epsilon^{a-\epsilon}\gamma_m(x)\gamma_n(a-x)\,dx
 +\int_a^1\gamma_m(x)\gamma_n(1+a-x)\,dx\\
 &+\frac{\gamma_n(a)}{m+1}\log^{m+1}\epsilon
  +\frac{\gamma_m(a)}{n+1}\log^{n+1}\epsilon\biggr\}.
\end{align*}
The use of cutoff singular coefficients at $a$ introduces errors at most
$O(\epsilon|\log\epsilon|^M)$ for a fixed $M$, hence no additional
constant. Integrating the explicitly subtracted logarithmic terms
evaluates this limit as \eqref{per:eq:Imn}. This agrees with convolution
defined by addition on the torus. Equations \eqref{per:eq:I00}--
\eqref{per:eq:I11} follow by restricting \eqref{per:eq:00}--
\eqref{per:eq:11} to $(0,1)$, where $\Delta$ is the constant $-1$.
\end{proof}

In terms of the ordinary digamma function, the first identity reads
\begin{align}
 2\gamma_1(a)+\zeta(2)={}&
 \int_0^a\left\{\psi(x)\psi(a-x)
       +\psi(a)\left(\frac1x+\frac1{a-x}\right)\right\}dx\notag\\
 &+\int_a^1\psi(x)\psi(1+a-x)\,dx-2\psi(a)\log a.
 \label{per:eq:digammaintegral}
\end{align}
No divergent integral occurs on the right: the braces are to be combined
before integration. For instance,
\[
 \mathcal I_{0,0}(1/2)
 =2\gamma_1-4\gamma\log2-2\log^22+\zeta(2),
\]
because $\gamma_1(1/2)=\gamma_1-2\gamma\log2-\log^22$.
The factor of the unshifted $\gamma_1$ in the latter identity is $1$,
as follows directly from $\zeta(s,1/2)=(2^s-1)\zeta(s)$.

\subsection{Every repeated Stieltjes primitive in one finite formula}

For a zero-mean periodic distribution define $D^{-q}$ by its Fourier
multiplier $(2\pi i n)^{-q}$ at $n\ne0$ and zero at $n=0$. This is
the unique repeated primitive having zero mean at each stage; it equals
convolution with $W_{1-q}$. Define the rational coefficients
\begin{equation}\label{per:eq:h}
 \prod_{r=1}^{q-1}\left(1-\frac wr\right)^{-1}
       =\sum_{j\ge0}h_j^{(q-1)}w^j,
 \qquad h_0^{(0)}=1,\quad h_j^{(0)}=0\ (j>0).
\end{equation}
They are complete homogeneous symmetric functions in
$1,1/2,\ldots,1/(q-1)$; equivalently their generating function is
\[
 \exp\left(\sum_{r\ge1}\frac{H_{q-1}^{(r)}}r w^r\right),
 \qquad H_N^{(r)}=\sum_{j=1}^N\frac1{j^r}.
\]
For example $h_1^{(N)}=H_N$ and
$h_2^{(N)}=(H_N^2+H_N^{(2)})/2$.

\begin{theorem}[All repeated primitives and their constants]
\label{per:thm:primitive}
For $q\ge1,n\ge0$, the zero-mean $q$-fold primitive of $T_n$ is
\begin{equation}\label{per:eq:primitive}
 D^{-q}T_n(x)=\frac{(-1)^{n+1}n!}{(q-1)!}
 \sum_{j=0}^{n+1}\frac{h_{n+1-j}^{(q-1)}}{j!}
                              \zeta^{[j]}(1-q,x),\quad 0<x<1.
\end{equation}
Here $\zeta^{[j]}$ means the $j$th derivative in the first argument.
The right side is periodized as a locally integrable function. For
$q\ge2$ it is continuous, with the regularity implied by its endpoint
terms $x^{q-1}\log^j x$.

If $q>k\ge0$ and $p=q-k$, then
\begin{align}
 D^{-q}P_k(x)&=
 \frac{\zeta'(1-p,x)+(H_{p-1}-H_k)\zeta(1-p,x)}{(p-1)!},
 \label{per:eq:Pprimitive}\\
 D^{-q}Q_k(x)&=
 \frac{\zeta'(1-p,x)+H_{p-1}\zeta(1-p,x)}{(p-1)!}.
 \label{per:eq:Qprimitive}
\end{align}
\end{theorem}

\begin{proof}
The normalized family gives, identically in $w$ near zero,
\begin{equation}\label{per:eq:primitivegenerator}
 D^{-q}Z_{1+w}
 =\Gamma(-w)W_{1-q+w}
 =\frac{Z_{1-q+w}}{(-w)_q}.
\end{equation}
Since
\[
 \frac1{(-w)_q}
 =-\frac1{(q-1)!w}\prod_{r=1}^{q-1}(1-w/r)^{-1},
\]
the coefficient $(-1)^nn![w^n]$ is exactly
\eqref{per:eq:primitive}. For $q\ge1$ the relevant Hurwitz jets have
order $1-q\le0$ and are locally integrable. Their means vanish, by the
integrable Hurwitz identity continued and differentiated below $1$,
so no hidden constants occur. For $q\ge2$, the shift formula isolates
the only nonanalytic endpoint term as a finite combination of
$x^{q-1}\log^j x$; it tends to zero and matches the finite endpoint
value at $1$.

For \eqref{per:eq:Pprimitive}, the Fourier coefficient after taking
$q$ primitives is $(2\pi i n)^{-p}(\gamma+L_n-H_k)$.
Differentiation of $Z_s=\Gamma(1-s)W_s$ at $s=1-p$ gives
\[
 \frac{Z'_{1-p}}{(p-1)!}
 =\partial_sW_s\big|_{s=1-p}-\psi(p)W_{1-p}.
\]
Use $\psi(p)=H_{p-1}-\gamma$. Adding the primitive of
$H_k\delta_0^{(k)}$ gives \eqref{per:eq:Qprimitive}.
\end{proof}

At $q=1$, \eqref{per:eq:primitive} becomes the familiar Stieltjes
primitive $(-1)^{n+1}\zeta^{[n+1]}(0,x)/(n+1)$. At every larger $q$
it gives all lower jets and all constants of integration explicitly.
For example,
\begin{align*}
 D^{-2}T_1(x)&=\tfrac12\zeta''(-1,x)+\zeta'(-1,x)+\zeta(-1,x),\\
 D^{-3}T_1(x)&=\tfrac14\zeta''(-2,x)+\tfrac34\zeta'(-2,x)
                                      +\tfrac78\zeta(-2,x).
\end{align*}
The family in \eqref{per:eq:Qprimitive} is the classical balanced
negative polygamma family. Equation \eqref{per:eq:Pprimitive} records
the correction required when one starts from spectral finite parts
instead of derivatives of periodic $\log\Gamma$.


\subsection{A convergent Gamma convolution as a normalization check}

One can remove the singularities by integrating. In particular, for
$g=\log\Gamma-\tfrac12\log(2\pi)$,
\begin{align}
 &\int_0^a g(x)g(a-x)\,dx
       +\int_a^1g(x)g(1+a-x)\,dx\notag\\
 &\hspace{10mm}=\zeta''(-1,a)+2\zeta'(-1,a)
                           +(2-\zeta(2))\zeta(-1,a),\quad 0<a\le1.
 \label{per:eq:gammaordinary}
\end{align}
Indeed $Dg=P_0$, so $g*g=D^{-2}(P_0*P_0)$; apply
\eqref{per:eq:00} and \eqref{per:eq:primitive}. The right side extends
continuously to the endpoints. This is a differentiated instance of
Sun's classical ordinary convolution, included to exhibit the agreement
of the singular calculus with convergent Gamma integrals.

% END input


% BEGIN sections/03_twisted.tex
\section{Nonintegral twists: Lerch jets, contact terms, and the half-twist}
\label{tw:sec:main}

The pinned convolution-algebra continuation explicitly asks for a
finite-part Stieltjes normal form for nonintegral frequency shifts, with
an all-index half-twist multiplication table that includes its
point-supported counterterms
\cite[``Twists, characters, and alternate endpoint normalizations'']{RepoConvolutionAlgebra}.
We give such a normalization and table here. The underlying Lerch
function is classical \cite[\S25.14]{DLMFLerch}; at the half-twist its
order jets are the modified Stieltjes constants of Hu and
Kim~\cite{HuKim2022}. An ordinary, nonsingular alternating Hurwitz
convolution is already given by Zhu, Hu, and
Kim~\cite[Theorem~1.13]{ZhuHuKim2026}. The additional content below is
the continuation through the singular orders, including the exact
contact terms, covariant primitives, and degeneration to the untwisted
normalization. In particular, twisting changes the special-function
jets and the convolution identity; the positive-integer zeta values in
the Gamma correction remain the same.

\subsection{The meromorphic family and its endpoint normalization}

Fix $0<\theta<1$, put $\omega=2\pi\theta$ and $z=e^{-i\omega}$, and
retain the Fourier and convolution conventions of
Section~\ref{per:sec:main}. Any other real nonintegral twist
$\theta=m+\vartheta$, $m\in\mathbb Z$, $0<\vartheta<1$, is obtained
by the periodic modulation $U(x)\mapsto e^{-2\pi imx}U(x)$.
Thus the chosen interval loses no real nonintegral shifts. Define
\begin{equation}\label{tw:eq:lambda}
 \lambda_{n,\theta}=2\pi i(n+\theta),\qquad
 L_{n,\theta}=\log(2\pi|n+\theta|)
                      +\frac{i\pi}{2}\operatorname{sgn}(n+\theta),
 \qquad n\in\mathbb Z.
\end{equation}
Every power of $\lambda_{n,\theta}$ uses this logarithm. There is no
zero frequency to remove. We write
\[
 D_\theta=D+2\pi i\theta,
 \qquad \widehat{D_\theta U}(n)=\lambda_{n,\theta}\widehat U(n).
\]
Here $D_\theta$ acts on ordinary periodic distributions, rather than on
functions satisfying a different boundary condition.
The nonperiodic factor $e^{-i\omega x}$ below is used only for
ordinary representatives on $0<x<1$ and for local endpoint Taylor
coefficients. It is not used as a global multiplier on periodic
distributions; the Fourier or cutoff definitions fix those extensions.
Explicitly, for $m\in\mathbb Z$ the extension to shifted intervals is
$W_s^{\theta+m}(x)=e^{-2\pi imx}W_s^\theta(x)$, with the same identity
for $Z_s$, $U_n$, and $P_k$ after their definitions below.

\begin{theorem}[Twisted Lerch family]\label{tw:thm:family}
The Fourier coefficients
\begin{equation}\label{tw:eq:W}
 \widehat{W_s^\theta}(n)=\lambda_{n,\theta}^{\,s-1},
 \qquad n\in\mathbb Z,
\end{equation}
define an entire distribution-valued function of $s$. The family
\begin{equation}\label{tw:eq:Z}
 Z_s^\theta=\Gamma(1-s)W_s^\theta
\end{equation}
is meromorphic, with simple poles precisely at $s=k+1$, $k\ge0$, and
\begin{equation}\label{tw:eq:residues}
 \operatorname{Res}_{s=k+1}Z_s^\theta
       =\frac{(-1)^{k+1}}{k!}D_\theta^k\delta_0.
\end{equation}
For $\operatorname{Re}s<1$, it is the locally integrable periodic
extension of
\begin{equation}\label{tw:eq:Lerch}
 f_s^\theta(x)=e^{-i\omega x}\Phi(z,s,x),\qquad 0<x<1.
\end{equation}
For every $s$, the restriction of the meromorphic distribution to
$(0,1)$ is the entire function of $s$ on the right of
\eqref{tw:eq:Lerch}. The convolution and differential identities are
\begin{align}
 W_s^\theta*W_t^\theta&=W_{s+t-1}^\theta,
 &W_1^\theta&=\delta_0,\label{tw:eq:group}\\
 D_\theta W_s^\theta&=W_{s+1}^\theta,
 &D_\theta Z_s^\theta&=-s Z_{s+1}^\theta.\label{tw:eq:derivative}
\end{align}
Consequently,
\begin{equation}\label{tw:eq:Zconv}
 Z_s^\theta*Z_t^\theta
 =\frac{\Gamma(1-s)\Gamma(1-t)}{\Gamma(2-s-t)}
                           Z_{s+t-1}^\theta
\end{equation}
as a meromorphic identity, interpreted through \eqref{tw:eq:group} at
intersections of polar divisors.
\end{theorem}

\begin{proof}
On compact sets of $s$, the coefficients in \eqref{tw:eq:W}, and each
fixed number of their order derivatives, have a common polynomial
growth bound in $|n|$. Pairing their Fourier series with the rapidly
decreasing Fourier coefficients of a smooth test function therefore
gives normal convergence. This proves entire dependence. The Fourier
multiplier of $D_\theta^k\delta_0$ is $\lambda_{n,\theta}^k$;
the residues of $\Gamma(1-s)$ prove \eqref{tw:eq:residues}. These
distributions are nonzero, so every displayed pole is present.

For completeness, the Lerch identification can be proved without
assuming a boundary-value Fourier formula. Initially let
$0<\operatorname{Re}s<1$ and $\epsilon>0$. Abel damping gives
\[
 f_{s,\epsilon}^\theta(x)
 =\sum_{m\ge0}e^{-(\epsilon+i\omega)(m+x)}(m+x)^{-s}.
\]
Absolute convergence permits unfolding its $n$th Fourier coefficient:
\begin{align*}
 \int_0^1 f_{s,\epsilon}^\theta(x)e^{-2\pi i n x}\,dx
 &=\int_0^\infty t^{-s}e^{-(\epsilon+\lambda_{n,\theta})t}\,dt\\
 &=\Gamma(1-s)(\epsilon+\lambda_{n,\theta})^{s-1}.
\end{align*}
Since $z\ne1$, the partial sums of $z^m$ are bounded. Summation by
parts shows that the tail starting at $m=1$ converges uniformly in
$x\in[0,1]$, and locally uniformly for $s$ in this strip, as
$\epsilon\downarrow0$. The remaining term is bounded by an integrable
multiple of $x^{-\operatorname{Re}s}$. Thus the limit can be taken in
$L^1(0,1)$ and in its Fourier coefficients, giving
\eqref{tw:eq:Lerch}. Repeated summation by parts, or finite differences
of $(m+x)^{-s}$, continues the Lerch tail holomorphically to every
$s$, locally uniformly when $x$ stays away from zero. Near zero its
only nonsmooth term is
\begin{equation}\label{tw:eq:localLerch}
 f_s^\theta(x)=e^{-i\omega x}x^{-s}
                +e^{-i\omega x}z\Phi(z,s,1+x).
\end{equation}
This proves $L^1$ holomorphy throughout $\operatorname{Re}s<1$ and
the asserted continuation away from the endpoint. The remaining
identities follow by multiplying the Fourier coefficients; the
Gamma recurrence gives the last identity of
\eqref{tw:eq:derivative}.
\end{proof}

The pointwise function in \eqref{tw:eq:Lerch} has no pole in $s$.
Every pole in \eqref{tw:eq:Z} is an endpoint distribution. This is
different from the untwisted Hurwitz family, whose pole at $s=1$
also has an ordinary constant contribution away from the endpoint.

For $n\ge0$ define ordinary Lerch jets and their periodic finite parts
by
\begin{align}
 c_n^\theta(x)&=(-1)^n
          \left.\partial_s^n f_s^\theta(x)\right|_{s=1},
               &&0<x<1,\label{tw:eq:ordinaryjets}\\
 Z_{1+u}^\theta&=-\frac{\delta_0}{u}
       +\sum_{n\ge0}\frac{(-1)^n}{n!}U_n^\theta u^n.
 \label{tw:eq:jets}
\end{align}
Set, also,
\begin{equation}\label{tw:eq:Pk}
 P_k^\theta=(-1)^{k+1}k!
               \operatorname{FP}_{s=k+1}Z_s^\theta,
 \qquad k\ge0.
\end{equation}
In particular $P_0^\theta=-U_0^\theta$. Off the endpoint it equals
$p_k^\theta(x)=(-1)^{k+1}k! f_{k+1}^\theta(x)$.

\begin{proposition}[Explicit singular normalization]
\label{tw:prop:cutoff}
For a smooth periodic test function $\phi$, put
$h(x)=e^{-i\omega x}\phi(x)$ near zero. Then
\begin{equation}\label{tw:eq:Ucutoff}
 \langle U_n^\theta,\phi\rangle
 =\lim_{\epsilon\downarrow0}\left\{
      \int_\epsilon^1 c_n^\theta(x)\phi(x)\,dx
        +\frac{\phi(0)}{n+1}\log^{n+1}\epsilon\right\}.
\end{equation}
If $F_k^\theta=\operatorname{FP}_{s=k+1}Z_s^\theta$ and $k\ge1$,
then
\begin{align}
 \langle F_k^\theta,\phi\rangle
 =\lim_{\epsilon\downarrow0}\biggl\{&
       \int_\epsilon^1 f_{k+1}^\theta(x)\phi(x)\,dx
       -\sum_{j=0}^{k-1}\frac{h^{(j)}(0)}{j!}
                         \frac{\epsilon^{j-k}}{k-j}
       +\frac{h^{(k)}(0)}{k!}\log\epsilon\biggr\}.
 \label{tw:eq:Pcutoff}
\end{align}
These normalizations retain the zero Fourier coefficient; no
mean-zero condition is imposed.
\end{proposition}

\begin{proof}
Write the second term in \eqref{tw:eq:localLerch} as
$r_s^\theta(x)$. After subtracting the degree-$N$ Taylor polynomial
of $h$, the continuation of the test-function pairing is
\begin{align}
 \langle Z_s^\theta,\phi\rangle={}&
 \int_0^1x^{-s}\left(h(x)-\sum_{j=0}^N
                     \frac{h^{(j)}(0)}{j!}x^j\right)dx
 +\sum_{j=0}^N\frac{h^{(j)}(0)}{j!(j+1-s)}\notag\\
 &+\int_0^1r_s^\theta(x)\phi(x)\,dx.
 \label{tw:eq:testsubtraction}
\end{align}
The first integral is holomorphic for $\operatorname{Re}s<N+2$,
and the last integral is entire. At $s=1+u$, expand
$x^{-1-u}=x^{-1}\sum_n(-1)^n u^n\log^n x/n!$ after the constant
subtraction. This gives \eqref{tw:eq:Ucutoff}. At $s=k+1$, take
the constant Laurent coefficient and integrate each subtracted
monomial explicitly; this gives \eqref{tw:eq:Pcutoff}.
The residue computed this way agrees with \eqref{tw:eq:residues},
because
\[
 \langle D_\theta^k\delta_0,\phi\rangle
       =(-1)^k h^{(k)}(0).
\]
This also fixes the signs of all the lower derivatives of $\delta_0$
inside $D_\theta^k\delta_0$.
\end{proof}

\subsection{All-index multiplication and covariant contact terms}

Let
\begin{equation}\label{tw:eq:B}
 B_\theta(u)=-uZ_{1+u}^\theta
      =\Gamma(1-u)W_{1+u}^\theta,
 \qquad
 G_\theta(u)=\sum_{n\ge0}\frac{(-1)^n}{n!}U_n^\theta u^n.
\end{equation}
Thus $B_\theta(u)=\delta_0-uG_\theta(u)$. All the following
equalities are analytic germs at zero in the algebra of periodic
distributions with its usual convolution identity $\delta_0$.

\begin{theorem}[Twisted Stieltjes normal form]
\label{tw:thm:normalform}
With $C(u,v)$ as in \eqref{per:eq:Bconv},
\begin{align}
 B_\theta(u)*B_\theta(v)&=C(u,v)B_\theta(u+v),
 \label{tw:eq:Bproduct}\\
 B_\theta(u)&=\exp_*\left(uP_0^\theta+
                \sum_{r\ge2}\frac{\zeta(r)}r u^r\delta_0\right).
 \label{tw:eq:Bellgenerator}
\end{align}
In particular,
\begin{equation}\label{tw:eq:Bell}
 U_n^\theta=\frac{(-1)^{n+1}}{n+1}
 \mathcal B_{n+1}^*
 \bigl(P_0^\theta,\zeta(2)\delta_0,2!\zeta(3)\delta_0,
                            \ldots,n!\zeta(n+1)\delta_0\bigr).
\end{equation}
An explicit all-index multiplication table is furnished by the
removable coefficient formula
\begin{align}
 G_\theta(u)*G_\theta(v)
 =\frac{1}{uv}\bigl\{&uG_\theta(u)+vG_\theta(v)
       -(u+v)C(u,v)G_\theta(u+v)\notag\\
       &+(C(u,v)-1)\delta_0\bigr\}.
 \label{tw:eq:table}
\end{align}
More precisely, $U_m^\theta*U_n^\theta$ is $(-1)^{m+n}m!n!$
times the coefficient of $u^mv^n$ on the right. This is a finite
combination of $\delta_0,U_0^\theta,\ldots,U_{m+n+1}^\theta$,
with coefficients polynomial in ordinary zeta values.
\end{theorem}

\begin{proof}
The first identity follows from \eqref{tw:eq:group}. At every
Fourier mode, including $n=0$,
\[
 \widehat{P_0^\theta}(n)=\gamma+L_{n,\theta},\qquad
 \widehat{B_\theta(u)}(n)
       =\exp\left((\gamma+L_{n,\theta})u+
                         \sum_{r\ge2}\frac{\zeta(r)}r u^r\right).
\]
This proves \eqref{tw:eq:Bellgenerator} and the Bell formula, with
normal convergence after test-function pairing. Substitute
$B_\theta(u)=\delta_0-uG_\theta(u)$ into
\eqref{tw:eq:Bproduct} and solve for
$G_\theta(u)*G_\theta(v)$ to obtain \eqref{tw:eq:table}. Its
numerator vanishes when either variable is zero, so the quotient
is removable. Extracting a fixed coefficient uses only finitely
many coefficients of the Gamma quotient and of $G_\theta$.
\end{proof}

For example, the full identities, including their contact masses, are
\begin{align}
 U_0^\theta*U_0^\theta
   &=2U_1^\theta-\zeta(2)\delta_0,\label{tw:eq:00}\\
 U_0^\theta*U_1^\theta
   &=\tfrac32U_2^\theta-\zeta(2)U_0^\theta
                                      +\zeta(3)\delta_0,\label{tw:eq:01}\\
 U_0^\theta*U_2^\theta
   &=\tfrac43U_3^\theta-2\zeta(2)U_1^\theta
           +2\zeta(3)U_0^\theta-2\zeta(4)\delta_0,\label{tw:eq:02}\\
 U_1^\theta*U_1^\theta
   &=U_3^\theta-2\zeta(2)U_1^\theta
           +2\zeta(3)U_0^\theta-\tfrac14\zeta(4)\delta_0.
 \label{tw:eq:11}
\end{align}
The presence of $\delta_0$, rather than $\delta_0-1$, matters even
when these formulas are eventually evaluated away from zero.

\begin{theorem}[All covariant contact corrections]
\label{tw:thm:contacts}
For $k\ge0$, with $H_0=0$,
\begin{align}
 \widehat{P_k^\theta}(n)
   &=\lambda_{n,\theta}^k
                 (\gamma+L_{n,\theta}-H_k),\label{tw:eq:PFourier}\\
 D_\theta P_k^\theta
   &=P_{k+1}^\theta+\frac{1}{k+1}D_\theta^{k+1}\delta_0,
 \label{tw:eq:DP}\\
 Q_k^\theta:=D_\theta^kP_0^\theta
   &=P_k^\theta+H_kD_\theta^k\delta_0.
 \label{tw:eq:Q}
\end{align}
Hence the all-index twisted polygamma product is
\begin{align}
 P_k^\theta*P_\ell^\theta
  =D_\theta^{k+\ell}\bigl\{
      2U_1^\theta+(H_k+H_\ell)U_0^\theta
                 +(H_kH_\ell-\zeta(2))\delta_0\bigr\}.
 \label{tw:eq:polygammaproduct}
\end{align}
More generally, let $R_{n,k}^\theta$ be the raw Hadamard finite part
of the ordinary function $D_\theta^kc_n^\theta$ on $(0,1)$, using
Taylor subtraction of its powers of $x$ and $\log x$ at zero and no
additional endpoint mass. Then
\begin{equation}\label{tw:eq:allcontacts}
 D_\theta^kU_n^\theta
 =R_{n,k}^\theta-h_{k,n}D_\theta^k\delta_0,
 \qquad
 h_{k,n}=(-1)^n n![u^{n+1}]
                         \prod_{j=1}^k\left(1+\frac{u}{j}\right).
\end{equation}
In particular $h_{k,0}=H_k$, and $h_{k,n}=0$ for $n\ge k$.
\end{theorem}

\begin{proof}
At $s=k+1+u$,
\[
 \Gamma(-k-u)
 =\frac{(-1)^{k+1}}{k!u}
       +\frac{(-1)^k}{k!}(H_k-\gamma)+O(u).
\]
Multiply this by
$\lambda_{n,\theta}^ke^{uL_{n,\theta}}$ and take the finite part.
This proves \eqref{tw:eq:PFourier}. The two derivative formulas
follow by multiplication by $\lambda_{n,\theta}$ and
$H_{k+1}-H_k=1/(k+1)$. Multiplying the two affine polynomials
$\gamma+L-H_k$ and $\gamma+L-H_\ell$ and using
\eqref{tw:eq:00} gives \eqref{tw:eq:polygammaproduct}.

For clarity, the raw finite part in \eqref{tw:eq:allcontacts} can
equivalently be defined by the coefficientwise constant cutoff term
of $D_\theta^kf_{1+u}^\theta$. The Mellin subtraction
\eqref{tw:eq:testsubtraction} identifies the regular Laurent
coefficients of $Z_{k+1+u}^\theta$ with those cutoff finite parts
of its order derivatives. Repeated application of
\eqref{tw:eq:derivative} gives
\[
 D_\theta^kZ_{1+u}^\theta
          =(-1)^k(1+u)_kZ_{k+1+u}^\theta.
\]
The contribution from the pole on the right is exactly
\[
 -\frac{1}{u}\prod_{j=1}^k\left(1+\frac{u}{j}\right)
                  D_\theta^k\delta_0.
\]
Its constant and higher coefficients supply all the difference
between derivative-compatible and raw finite parts. Equating
the coefficients of $(-1)^nu^n/n!$ proves
\eqref{tw:eq:allcontacts}.
\end{proof}

For example, $D_\theta^2\delta_0=\delta_0''+
4\pi i\theta\delta_0'-(2\pi\theta)^2\delta_0$. Thus
\eqref{tw:eq:Q} specifies the lower point-supported terms as well
as the highest derivative. Writing only $H_k\delta_0^{(k)}$ would
give the wrong normalization for a nonzero twist.

\subsection{Rational twists and ordinary special functions}

\begin{proposition}[Finite Hurwitz reduction]\label{tw:prop:rational}
Let $\theta=p/q$, where $1\le p<q$ are relatively prime, and put
$z=e^{-2\pi ip/q}$. Then
\begin{equation}\label{tw:eq:rationalLerch}
 f_s^{p/q}(x)=\frac{e^{-2\pi ipx/q}}{q^s}
          \sum_{r=0}^{q-1}z^r\zeta\left(s,\frac{x+r}{q}\right),
 \qquad 0<x<1.
\end{equation}
The sum is entire in $s$: its Hurwitz residues cancel. Consequently,
\begin{align}
 c_n^{p/q}(x)
 &=\frac{e^{-2\pi ipx/q}}q
       \sum_{r=0}^{q-1}z^r\sum_{j=0}^n\binom nj
        (\log q)^{n-j}\gamma_j\left(\frac{x+r}{q}\right),
 \label{tw:eq:rationalStieltjes}\\
 p_k^{p/q}(x)
 &=\frac{e^{-2\pi ipx/q}}{q^{k+1}}
       \sum_{r=0}^{q-1}z^r\psi^{(k)}\left(\frac{x+r}{q}\right),
       \qquad k\ge0.\label{tw:eq:rationalPolygamma}
\end{align}
These ordinary functions, extended by
\eqref{tw:eq:Ucutoff} and \eqref{tw:eq:Pcutoff}, give the
finite-part distributions in Theorems~\ref{tw:thm:normalform}
and~\ref{tw:thm:contacts}.
\end{proposition}

\begin{proof}
For $\operatorname{Re}s>1$, split the defining Lerch series into
its residue classes modulo $q$. This gives
\eqref{tw:eq:rationalLerch}; analytic continuation proves it
everywhere. Since $\sum_r z^r=0$, the common Hurwitz pole vanishes
before the factor $q^{-s}$ is expanded. Differentiating the resulting
Taylor expansion at $s=1$ gives
\eqref{tw:eq:rationalStieltjes}. At positive integers use the
Hurwitz--polygamma identity; its limiting case at $s=1$ uses
$\gamma_0=-\psi$ and the same pole cancellation. This proves
\eqref{tw:eq:rationalPolygamma} also for $k=0$.
\end{proof}

Thus additive rational twists require finite colored combinations of
Hurwitz jets, while the coefficients governing their convolution
remain the ordinary $\zeta(r)$ from $\log\Gamma(1-u)$.
No replacement by Dirichlet eta values is needed in
\eqref{tw:eq:Bell} or \eqref{tw:eq:table}.

\subsection{The half-twist and convergent alternating identities}

At $\theta=\tfrac12$, set
\begin{align}
 \zeta_E(s,x)&=\Phi(-1,s,x)
       =2^{-s}\left\{\zeta\left(s,\frac{x}{2}\right)
                  -\zeta\left(s,\frac{x+1}{2}\right)\right\},
 \label{tw:eq:Eulerzeta}\\
 \widetilde\gamma_n(x)&=(-1)^n
                   \left.\partial_s^n\zeta_E(s,x)\right|_{s=1}.
 \label{tw:eq:modified}
\end{align}
These modified Stieltjes constants are the established functions of
Hu and Kim~\cite{HuKim2022}. The explicit conversion is
\begin{equation}\label{tw:eq:halfStieltjes}
 \widetilde\gamma_n(x)
 =\frac12\sum_{j=0}^n\binom nj(\log2)^{n-j}
       \left\{\gamma_j\left(\frac{x}{2}\right)
              -\gamma_j\left(\frac{x+1}{2}\right)\right\},
 \qquad c_n^{1/2}(x)=e^{-\pi ix}\widetilde\gamma_n(x).
\end{equation}
In particular, the Nielsen beta function is
\begin{equation}\label{tw:eq:beta}
 \beta(x)=\widetilde\gamma_0(x)
      =\frac12\left\{\psi\left(\frac{x+1}{2}\right)
                     -\psi\left(\frac{x}{2}\right)\right\}.
\end{equation}
The elementary shift relation
\begin{equation}\label{tw:eq:modifiedshift}
 \widetilde\gamma_n(x)+\widetilde\gamma_n(x+1)
                                      =\frac{\log^n x}{x}
\end{equation}
fixes its endpoint singularity, including every logarithmic order.

\begin{theorem}[All-index alternating finite-part convolution]
\label{tw:thm:alternating}
For $0<a<1$ and $m,n\ge0$, the following expression consists of
absolutely convergent integrals:
\begin{align}
 \mathcal A_{mn}(a)={}&
 \int_0^a\biggl\{
      \widetilde\gamma_m(x)\widetilde\gamma_n(a-x)
       -\widetilde\gamma_n(a)\frac{\log^m x}{x}
       -\widetilde\gamma_m(a)\frac{\log^n(a-x)}{a-x}
                 \biggr\}\,dx\notag\\
 &-\int_a^1\widetilde\gamma_m(x)
                        \widetilde\gamma_n(1+a-x)\,dx\notag\\
 &+\frac{\widetilde\gamma_n(a)}{m+1}\log^{m+1}a
       +\frac{\widetilde\gamma_m(a)}{n+1}\log^{n+1}a.
 \label{tw:eq:Amn}
\end{align}
It equals $e^{\pi ia}(U_m^{1/2}*U_n^{1/2})(a)$, where the
convolution is smooth on $(0,1)$. Every $\mathcal A_{mn}$ is
therefore the finite linear combination of
$\widetilde\gamma_0,\ldots,\widetilde\gamma_{m+n+1}$ obtained
from \eqref{tw:eq:table} by discarding its $\delta_0$ term.
In particular,
\begin{align}
 \mathcal A_{00}(a)&=2\widetilde\gamma_1(a),
       \label{tw:eq:A00}\\
 \mathcal A_{01}(a)&=\tfrac32\widetilde\gamma_2(a)
                          -\zeta(2)\widetilde\gamma_0(a),
       \label{tw:eq:A01}\\
 \mathcal A_{11}(a)&=\widetilde\gamma_3(a)
             -2\zeta(2)\widetilde\gamma_1(a)
             +2\zeta(3)\widetilde\gamma_0(a).
       \label{tw:eq:A11}
\end{align}
\end{theorem}

\begin{proof}
For integrable orders, convolution on the circle splits into
$0<x<a$ and $a<x<1$. The phases of the two factors have product
$e^{-\pi ia}$ in the first interval and
$e^{-\pi i(1+a)}=-e^{-\pi ia}$ in the second. This is the source
of the alternating sign in \eqref{tw:eq:Amn}, and agrees with
the nonsingular alternating convolution of
\cite[Theorem~1.13]{ZhuHuKim2026}.

To continue to the jets, use the normalized cutoff distributions in
\eqref{tw:eq:Ucutoff}. On a compact subinterval of $0<a<1$, the
two endpoint singularities in the first integral occur at distinct
points, $x=0$ and $x=a$. Their only divergent terms are the two
explicit subtractions in \eqref{tw:eq:Amn}. For instance, at zero,
\eqref{tw:eq:modifiedshift} writes the first factor as
$\log^m x/x$ plus a smooth function, while
$\widetilde\gamma_n(a-x)-\widetilde\gamma_n(a)=O(x)$.
The subtracted integrand is therefore
$O(1+|\log x|^m)$ there, and has the analogous integrable behavior
at $a$. The wrapped integral is smooth. Integration of the two
subtracted singular monomials gives exactly the two finite
$\log a$ terms shown. This verifies the cutoff expression and
permits all parameter coefficients to be taken under the
subtracted integrals. Equivalently it is the restriction of the
meromorphic convolution identity \eqref{tw:eq:Zconv}. The identities
\eqref{tw:eq:A00}--\eqref{tw:eq:A11} now follow from
\eqref{tw:eq:00}, \eqref{tw:eq:01}, and \eqref{tw:eq:11}.
\end{proof}

For ordinary digamma functions this gives the concrete identity
\begin{align}
 &\int_0^a\left\{\beta(x)\beta(a-x)
               -\beta(a)\left(\frac1x+\frac1{a-x}\right)\right\}dx
       -\int_a^1\beta(x)\beta(1+a-x)\,dx
       +2\beta(a)\log a\notag\\
 &\hspace{1cm}=\gamma_1\left(\frac a2\right)
          -\gamma_1\left(\frac{a+1}{2}\right)
                         +2\log2\,\beta(a).
 \label{tw:eq:betaIntegral}
\end{align}
The full distributional identity behind it still contains
$-\zeta(2)\delta_0$. Omitting that mass would preserve
\eqref{tw:eq:betaIntegral} on $0<a<1$ but would destroy the
Fourier coefficients and every identity involving endpoint tests.

\begin{corollary}[A quarter-Gamma evaluation]
\label{tw:cor:quarterGamma}
At $a=\tfrac12$, the absolutely convergent integral in
\eqref{tw:eq:betaIntegral} has the exact value
\begin{equation}\label{tw:eq:quarterGamma}
 \mathcal A_{00}\left(\tfrac12\right)
 =\gamma_1\left(\tfrac14\right)-\gamma_1\left(\tfrac34\right)
                                      +\pi\log2
 =\pi\left\{4\log\Gamma\left(\tfrac14\right)
                          -\gamma-3\log(2\pi)\right\}.
\end{equation}
\end{corollary}

\begin{proof}
Distinguish the Dirichlet beta function
$\beta_{\mathrm D}(s)=\sum_{r\ge0}(-1)^r(2r+1)^{-s}$
from the Nielsen function $\beta(x)$ in \eqref{tw:eq:beta}.
The identity $\zeta_E(s,\tfrac12)=2^s\beta_{\mathrm D}(s)$ gives
\[
 2\widetilde\gamma_1\left(\tfrac12\right)
                    =-\pi\log2-4\beta_{\mathrm D}'(1),
\]
since $\beta_{\mathrm D}(1)=\pi/4$.
The completed Dirichlet functional equation
\cite[\S25.15(i)]{DLMFDirichlet} is, for this character,
\[
 \Lambda(s)=\left(\frac4\pi\right)^{(s+1)/2}
             \Gamma\left(\frac{s+1}{2}\right)\beta_{\mathrm D}(s),
 \qquad \Lambda(s)=\Lambda(1-s).
\]
Also
\[
 \beta_{\mathrm D}(s)
  =4^{-s}\left\{\zeta\left(s,\tfrac14\right)
                         -\zeta\left(s,\tfrac34\right)\right\},
 \quad
 \beta_{\mathrm D}(0)=\tfrac12,\quad
 \beta_{\mathrm D}'(0)
  =\log\Gamma\left(\tfrac14\right)
                  -\log\Gamma\left(\tfrac34\right)-\log2.
\]
Logarithmic differentiation of the completed functional equation
at $0$ and $1$, followed by the Gamma reflection formula, yields
\[
 \beta_{\mathrm D}'(1)
   =\frac\pi4\left\{\gamma+2\log2+3\log\pi
                             -4\log\Gamma\left(\tfrac14\right)\right\}.
\]
Substitution proves \eqref{tw:eq:quarterGamma}; the first equality
also uses $\beta(\tfrac12)=\pi/2$ in
\eqref{tw:eq:betaIntegral}.
\end{proof}

\subsection{Reflection at the half-twist}

For the half-twist, the map $n\mapsto-n-1$ reverses the shifted
frequency $n+\tfrac12$. Its periodic realization is
\begin{equation}\label{tw:eq:J}
 \mathcal J U(x)=e^{-2\pi ix}U(-x),\qquad
 \widehat{\mathcal J U}(n)=\widehat U(-n-1).
\end{equation}
This is an involutive convolution automorphism with
$\mathcal J\delta_0=\delta_0$ and
$D_{1/2}\mathcal J=-\mathcal JD_{1/2}$. Ordinary reflection alone
does not reverse the half-shifted frequencies. Also define
\[
 \widehat{\mathcal H_{1/2}U}(n)
       =-i\operatorname{sgn}\left(n+\tfrac12\right)\widehat U(n),
 \qquad n\in\mathbb Z.
\]
Here $\mathcal H_{1/2}^2=-\mathrm{Id}$ on all periodic
distributions, because every shifted frequency is nonzero.

\begin{corollary}[Reflected half-twist table]
\label{tw:cor:reflection}
Writing $U_n=U_n^{1/2}$ and $B=B_{1/2}$ in this corollary,
\begin{equation}\label{tw:eq:reflectedB}
 B(u)*\mathcal J B(v)
  =C(u,v)\bigl(\cos(\pi v)\,\mathrm{Id}
                   +\sin(\pi v)\,\mathcal H_{1/2}\bigr)B(u+v).
\end{equation}
Thus the coefficient prescription gives an all-index reflected
table as well. Two explicit instances are
\begin{align}
 U_0*\mathcal J U_1
  &=\tfrac12 U_2+\mathcal J U_2
        +\zeta(2)(U_0+\mathcal J U_0)+\zeta(3)\delta_0,
 \label{tw:eq:reflected01}\\
 U_1*\mathcal J U_1
  &=\tfrac12(U_3+\mathcal J U_3)
        +2\zeta(2)(U_1+\mathcal J U_1)
        +\zeta(3)(U_0+\mathcal J U_0)
        +\tfrac72\zeta(4)\delta_0.
 \label{tw:eq:reflected11}
\end{align}
For the polygamma finite parts,
\begin{align}
 P_k^{1/2}*\mathcal J P_\ell^{1/2}
 =(-1)^\ell D_{1/2}^{k+\ell}\bigl\{
     U_1+\mathcal J U_1+H_\ell U_0+H_k\mathcal J U_0
               +(H_kH_\ell+2\zeta(2))\delta_0\bigr\}.
 \label{tw:eq:reflectedP}
\end{align}
\end{corollary}

\begin{proof}
At mode $n$ put $\sigma=\operatorname{sgn}(n+\tfrac12)$. The
logarithms of $\lambda_{n,1/2}$ and $-\lambda_{n,1/2}$ differ by
$i\pi\sigma$. Dividing the Fourier coefficient of the left side
of \eqref{tw:eq:reflectedB} by that of $B(u+v)$ gives
$C(u,v)e^{-i\pi v\sigma}$, which is exactly the multiplier on
the right. Coefficient extraction, with
$\zeta(2)^2=\tfrac52\zeta(4)$, proves the two displayed
instances. Multiplication of
$\gamma+L_{n,1/2}-H_k$ and
$\gamma+L_{-n-1,1/2}-H_\ell$ proves
\eqref{tw:eq:reflectedP}; the factor $(-1)^\ell$ comes from
$\lambda_{-n-1,1/2}^{\ell}=(-\lambda_{n,1/2})^\ell$.
\end{proof}

The modulation in \eqref{tw:eq:J} is periodic specifically because
$2\theta=1$. For a general twist, frequency reversal connects
$\theta$ with $1-\theta$; the same fixed-twist involution should
not be asserted without changing the space or the twist.

\subsection{Unique covariant primitives}

Because $\lambda_{n,\theta}\ne0$ for every $n$, the operator
$D_\theta$ has a unique inverse on periodic distributions:
\[
 \widehat{D_\theta^{-q}U}(n)
       =\lambda_{n,\theta}^{-q}\widehat U(n),\qquad q\ge1.
\]
This removes the choice of an integration constant that occurs
for ordinary periodic primitives.

\begin{theorem}[All-order covariant antiderivatives]
\label{tw:thm:primitives}
Define the complete homogeneous reciprocal-harmonic coefficients by
\begin{equation}\label{tw:eq:hcomplete}
 \sum_{j\ge0}h_j^{(q-1)}w^j
       =\prod_{r=1}^{q-1}\left(1-\frac wr\right)^{-1};
 \qquad h_0^{(0)}=1,\quad h_j^{(0)}=0\ (j>0).
\end{equation}
Then for $q\ge1$ and $n\ge0$,
\begin{equation}\label{tw:eq:primitive}
 D_\theta^{-q}U_n^\theta
 =\frac{(-1)^{n+1}n!}{(q-1)!}
       \sum_{j=0}^{n+1}\frac{h_{n+1-j}^{(q-1)}}{j!}
            \left.\partial_s^j Z_s^\theta\right|_{s=1-q}.
\end{equation}
Its ordinary value on $(0,1)$ is obtained by replacing
$Z_s^\theta$ by $f_s^\theta$. For $q\ge2$ this is a periodic
$C^{q-2}$ function. For $q>k\ge0$, setting $p=q-k$ gives
\begin{align}
 D_\theta^{-q}P_k^\theta
  &=\frac{
       \left.\partial_s Z_s^\theta\right|_{s=1-p}
           +(H_{p-1}-H_k)Z_{1-p}^\theta}{(p-1)!},
       \label{tw:eq:Pprimitive}\\
 D_\theta^{-q}Q_k^\theta
  &=\frac{
       \left.\partial_s Z_s^\theta\right|_{s=1-p}
           +H_{p-1}Z_{1-p}^\theta}{(p-1)!}.
       \label{tw:eq:Qprimitive}
\end{align}
At the half-twist, the first primitive has the explicit Gamma form
\begin{equation}\label{tw:eq:halfGamma}
 D_{1/2}^{-1}U_0^{1/2}(x)
  =e^{-\pi ix}\left\{
      \log\Gamma\left(\frac{x+1}{2}\right)
      -\log\Gamma\left(\frac{x}{2}\right)+\frac12\log2\right\},
 \qquad 0<x<1.
\end{equation}
The displayed additive term is fixed by the periodic distributional
normalization.
\end{theorem}

\begin{proof}
The spectral definition and the Gamma recurrence give
\[
 D_\theta^{-q}Z_{1+u}^\theta
       =\frac{Z_{1-q+u}^\theta}{(-u)_q},\qquad
 \frac1{(-u)_q}
       =-\frac1{(q-1)!u}\sum_{j\ge0}h_j^{(q-1)}u^j.
\]
The family in the numerator is holomorphic at $u=0$. Comparing
the coefficients of $(-1)^nu^n/n!$ proves
\eqref{tw:eq:primitive}. Its Fourier coefficients are
$O(|n|^{-q}(1+\log|n|)^{n_0+1})$ when the fixed Stieltjes
index is denoted by $n_0$. Therefore the Fourier series and its
first $q-2$ ordinary derivatives converge absolutely for $q\ge2$.
The derivative formulas for $P_k^\theta$ and $Q_k^\theta$ follow
directly from \eqref{tw:eq:PFourier} and
$\psi(p)=H_{p-1}-\gamma$.

For \eqref{tw:eq:halfGamma}, the case $q=1,n=0$ is
$-\partial_s Z_s^{1/2}|_{s=0}$. Use
\eqref{tw:eq:Eulerzeta}, $\zeta(0,a)=\tfrac12-a$, and
$\zeta'(0,a)=\log\Gamma(a)-\tfrac12\log(2\pi)$.
The difference of the two zeta values at zero is $\tfrac12$,
which supplies the term $\tfrac12\log2$ after differentiation.
\end{proof}

\subsection{Degeneration to the untwisted finite parts}

The zero Fourier coefficient becomes singular as $\theta\downarrow0$.
Deleting it exactly gives a continuous transition to the untwisted
normalization, including all the Stieltjes jets.

\begin{theorem}[Exact zero-mode renormalization]
\label{tw:thm:zerolimit}
For every $\rho<1$,
\begin{equation}\label{tw:eq:Blimit}
 B_\theta(u)-\Gamma(1-u)(2\pi i\theta)^u\,1
                    \longrightarrow B(u)
 \qquad (\theta\downarrow0)
\end{equation}
locally uniformly for $|u|<\rho$ after pairing with any smooth
periodic test function, and with every fixed number of derivatives
in $u$. Here $1$ is the constant distribution, which has only the
zero Fourier coefficient. Put
\begin{equation}\label{tw:eq:meanjet}
 \mu_n(\theta)=\frac{(-1)^{n+1}}{n+1}
 \mathcal B_{n+1}\bigl(
    \gamma+\log(2\pi i\theta),\zeta(2),2!\zeta(3),
                              \ldots,n!\zeta(n+1)\bigr).
\end{equation}
Then $\mu_n(\theta)=\widehat{U_n^\theta}(0)$ and
\begin{equation}\label{tw:eq:Ulimit}
 U_n^\theta-\mu_n(\theta)\,1\longrightarrow T_n
                \qquad\hbox{in }\mathcal D'(\mathbb R/\mathbb Z).
\end{equation}
In particular,
\[
 \mu_0(\theta)=-\gamma-\log(2\pi i\theta),\qquad
 \mu_1(\theta)=\frac{(\gamma+\log(2\pi i\theta))^2+\zeta(2)}2.
\]
\end{theorem}

\begin{proof}
The subtracted quantity in \eqref{tw:eq:Blimit} is exactly the
$n=0$ Fourier coefficient. For each $n\ne0$, the remaining
coefficient tends to
$\Gamma(1-u)(2\pi in)^u=\widehat{B(u)}(n)$ with the previously
fixed branch. Uniformly for $0<\theta\le\tfrac12$ and $u$ on a
compact subdisk, these coefficients and their fixed order
derivatives are bounded by a common polynomial in $|n|$, because
$|n+\theta|$ is comparable with $|n|$. Rapid decrease of the
test-function Fourier coefficients proves normal convergence and
allows termwise order differentiation. Taking the coefficient of
$(-1)^{n+1}u^{n+1}/n!$ gives \eqref{tw:eq:Ulimit}; the scalar
Gamma expansion gives \eqref{tw:eq:meanjet}.
\end{proof}

At $u=0$, the left side of \eqref{tw:eq:Blimit} is already
$\delta_0-1=\Delta$. Thus the change of convolution identity is
caused exactly by the disappearing zero mode. The same observation
explains why the interior alternating identity
\eqref{tw:eq:A00} has no added constant, whereas the untwisted
digamma convolution has the constant contributed by
$-\zeta(2)\Delta$.

% END input


% BEGIN sections/04_harmonic_laurent.tex
\section{Shifted harmonic sums, beta derivatives, and exact Laurent data}
\label{harm:sec:main}

The all-order Stieltjes differentiation formulas in the manuscript have
elementary symmetric harmonic coefficients.  Those coefficients can also
be put inside a Dirichlet series.  This gives a second exact calculus:
positive integer values are beta derivatives, whereas all principal parts
and finite parts at nonpositive integers are finite polynomials.  The
dependence on an arbitrary shift is retained throughout.

There are important antecedents.  Young's parametric hyperharmonic Euler
sums include the beta-derivative evaluation below
\cite{Young2026}; his work on height-one multiple zeta functions supplies
the unshifted generating-function framework \cite{Young2024}, and his
earlier paper treats their unshifted Laurent expansions \cite{Young2023}.  Harmonic
Stieltjes constants at depth one have a substantial separate literature
\cite{HarmonicStieltjes2023}.  Accordingly, neither the beta method nor the
unshifted height-one theory is asserted to be new.  The main continuation
developed here is the all-depth coefficient rule for outer-shifted
Laurent data at nonpositive integers, together with its polynomial shift
calculus and the residue term required when finite parts are
differentiated.  These formulas were not located in the inspected
repository sources; no claim of priority over all literature is made.

\subsection{Definitions and a single two-variable family}

Put $A=a+1$ and initially assume $\Re a>-1$.  Write
\begin{equation}\label{harm:eq:elementary}
 e_r(N)=[u^r]\prod_{j=1}^{N}\left(1+\frac uj\right),
 \qquad e_0(N)=1,
 \qquad e_r(N)=0\quad(r>N).
\end{equation}
Thus $e_r(N)=H_N(\{1\}^r)$ is the strict multiple harmonic sum, and
\begin{align*}
 e_1(N)&=H_N,\\
 e_2(N)&=\frac{H_N^2-H_N^{(2)}}2,\\
 e_3(N)&=\frac{H_N^3-3H_NH_N^{(2)}+2H_N^{(3)}}6.
\end{align*}
The main family is
\begin{equation}\label{harm:eq:E}
 E_r(s;a)=\sum_{n=1}^{\infty}\frac{e_r(n-1)}{(n+a)^s},
 \qquad r\geq0,\quad \Re s>1.
\end{equation}
The power uses the principal logarithm; $n+a$ lies in the right half-plane.
For $a=0$, this is $\zeta(s,\{1\}^r)$, and for $r=0$ it is
$\zeta(s,A)$.  The shift affects only the outer denominator.  In
particular, for general $a$ this is not the harmonic Hurwitz function
in which every inner denominator is shifted as well.

The bound $e_r(n-1)\leq H_{n-1}^r/r!$ proves normal convergence of
\eqref{harm:eq:E} on compact subsets of the displayed domain.  To keep
all depths together, define
\begin{equation}\label{harm:eq:F}
 F(s,u;a)=\sum_{n=1}^{\infty}
 \frac{\Gamma(n+u)}{\Gamma(n)\Gamma(1+u)}\frac1{(n+a)^s}.
\end{equation}
For $u$ near zero and $\Re s>1+|u|$, the defining series is normally
convergent.  The quotient in the numerator is the polynomial
$\prod_{j=1}^{n-1}(1+u/j)$, so
\begin{equation}\label{harm:eq:coeff}
 F(s,u;a)=\sum_{r\geq0}E_r(s;a)u^r.
\end{equation}
Equivalently, $F$ is the Barnes--Hurwitz zeta function with complex
order $1+u$, in the convention used in \cite{Young2024}.

\begin{proposition}[Mellin representations]\label{harm:prop:mellin}
In a common domain of absolute convergence,
\begin{align}
 \Gamma(s)F(s,u;a)
 &=\int_0^{\infty}t^{s-1}e^{-At}
                  (1-e^{-t})^{-1-u}\,dt,\label{harm:eq:mellinF}\\
 E_r(s;a)
 &=\frac1{r!\Gamma(s)}\int_0^1
       \frac{x^a(-\log x)^{s-1}[-\log(1-x)]^r}{1-x}\,dx.
       \label{harm:eq:mellinE}
\end{align}
For fixed $r$, the second integral converges for $\Re s>1$ and
$\Re a>-1$.
\end{proposition}
\begin{proof}
The binomial series gives
\[
 \sum_{n\geq1}\frac{\Gamma(n+u)}{\Gamma(n)\Gamma(1+u)}x^n
       =\frac{x}{(1-x)^{1+u}}.
\]
Multiply by $x^{a-1}(-\log x)^{s-1}$ and integrate.  Absolute
convergence permits termwise integration, and the integral of the
$n$th term is $\Gamma(s)(n+a)^{-s}$.  The substitution $x=e^{-t}$
gives \eqref{harm:eq:mellinF}; coefficient extraction in $u$ gives
\eqref{harm:eq:mellinE}.  At $x=1$, its absolute integrand is bounded
by a constant times
$(1-x)^{\Re s-2}(1+|\log(1-x)|^r)$, which is integrable for
$\Re s>1$.  At zero, $\Re a>-1$ suffices.
\end{proof}

This also gives a direct polylogarithmic interpretation, since
\[
 \operatorname{Li}_{\{1\}^{r+1}}(x)
 =\frac{[-\log(1-x)]^{r+1}}{(r+1)!},\qquad
 \frac{d}{dx}\operatorname{Li}_{\{1\}^{r+1}}(x)
 =\frac{[-\log(1-x)]^r}{r!(1-x)}.
\]
Thus \eqref{harm:eq:mellinE} is the Mellin transform of a derivative
of a multiple polylogarithm of arbitrary depth.

\subsection{Positive integer values and their antiderivatives}

Define generalized harmonic functions by
\begin{equation}\label{harm:eq:Hcomplex}
 H_a^{(1)}=\psi(A)+\gamma,\qquad
 H_a^{(k)}=\zeta(k)-\zeta(k,A)\quad(k\geq2).
\end{equation}
They agree with the usual finite harmonic numbers at nonnegative
integer $a$.  Introduce
\begin{equation}\label{harm:eq:R}
 R_a(u)=\frac{\Gamma(1-u)\Gamma(A)}{\Gamma(A-u)}
       =\exp\left(\sum_{k\geq1}\frac{H_a^{(k)}}k u^k\right)
       =\sum_{d\geq0}h_d(a)u^d.
\end{equation}
Here $h_d(a)$ is a complete symmetric harmonic polynomial:
\begin{equation}\label{harm:eq:hrec}
 h_0(a)=1,\qquad
 d\,h_d(a)=\sum_{k=1}^d H_a^{(k)}h_{d-k}(a).
\end{equation}
For example,
\[
 h_1=H_a^{(1)},\quad
 h_2=\frac{(H_a^{(1)})^2+H_a^{(2)}}2,\quad
 h_3=\frac{(H_a^{(1)})^3+3H_a^{(1)}H_a^{(2)}+2H_a^{(3)}}6.
\]
The plus signs distinguish these polynomials from the elementary
coefficients $e_r(N)$ in the summand.

\begin{theorem}[Subtracted beta identity and all positive integer values]
\label{harm:thm:beta}
For $r\geq0$ and $\Re a>-1$,
\begin{equation}\label{harm:eq:primitive}
 \boxed{\quad
 h_{r+1}(a)=\sum_{n\geq1}e_r(n-1)
           \left(\frac1n-\frac1{n+a}\right).
 \quad}
\end{equation}
For every integer $p\geq2$,
\begin{equation}\label{harm:eq:positive}
 \boxed{\quad
 E_r(p;a)=\frac{(-1)^p}{(p-1)!}
               \frac{d^{p-1}}{da^{p-1}}h_{r+1}(a).
 \quad}
\end{equation}
Equivalently, as an analytic germ at $(x,u)=(0,0)$,
\begin{equation}\label{harm:eq:betaG}
 \sum_{p\geq2}\sum_{r\geq0}E_r(p;a)x^{p-1}u^r
       =B(A-x,-u)-B(A,-u).
\end{equation}
The poles of the individual beta factors at $u=0$ cancel in this
difference.
\end{theorem}
\begin{proof}
Initially take $\Re u<0$.  The beta integral and
\eqref{harm:eq:mellinF} at $s=1$ give
\[
 F(1,u;a)=B(A,-u)=-\frac{R_a(u)}u,
 \qquad F(1,u;0)=-\frac1u.
\]
The difference of their series is normally convergent also for $u$
near zero: its $n$th term is $O(n^{|u|-2})$ on compact $a$-sets.
Consequently
\[
 \sum_{n\geq1}\prod_{j<n}(1+u/j)
       \left(\frac1n-\frac1{n+a}\right)
          =\frac{R_a(u)-1}{u}
\]
extends holomorphically through $u=0$.  Compare coefficients to
obtain \eqref{harm:eq:primitive}.  Normal convergence permits
$p-1$ derivatives in $a$, giving \eqref{harm:eq:positive}.

For an independent verification of the bivariate formula, sum
\eqref{harm:eq:mellinE} first over $p\geq2$ and then over $r$.
The result is the convergent integral
\[
 \int_0^1x_0^{A-1}(x_0^{-x}-1)(1-x_0)^{-1-u}\,dx_0.
\]
For $|x|$ sufficiently small and $|u|<1/2$, it is integrable at
both endpoints.  In the subdomain $\Re u<0$ it is a difference of
two beta integrals.  Analytic continuation across $u=0$ proves
\eqref{harm:eq:betaG} without separating divergent integrals there.
Finally, Taylor expansion of $\log\Gamma$ proves \eqref{harm:eq:R},
and logarithmic differentiation gives \eqref{harm:eq:hrec}.
\end{proof}

Formula \eqref{harm:eq:positive} is a specialization and explicit
coefficient implementation of the established parametric beta theory
\cite{Young2026}.  Its useful feature for this manuscript is that it
simultaneously evaluates infinitely many nonlinear harmonic combinations
and gives their antiderivatives.  In particular,
$\int E_r(2;a)\,da=h_{r+1}(a)+C$.

Put $Z_k=\zeta(k,A)$ and $H=H_a^{(1)}$.  Since
$\partial_a H_a^{(k)}=kZ_{k+1}$, the first outer exponent has the
especially short all-depth expression
\begin{equation}\label{harm:eq:second}
 E_r(2;a)=\sum_{k=1}^{r+1} Z_{k+1}h_{r+1-k}(a).
\end{equation}
For example,
\begin{align}
 E_1(2;a)&=HZ_2+Z_3,\label{harm:eq:examples1}\\
 E_2(2;a)&=h_2(a)Z_2+HZ_3+Z_4,\\
 E_1(3;a)&=HZ_3+\frac32Z_4-\frac12Z_2^2,\\
 E_2(3;a)&=h_2(a)Z_3+\frac32HZ_4+2Z_5
                   -\frac12HZ_2^2-Z_2Z_3.\label{harm:eq:examples4}
\end{align}
At $a=0$, \eqref{harm:eq:second} recovers
$\zeta(2,\{1\}^r)=\zeta(r+2)$; the next two examples give
$\zeta(3,1)=\zeta(4)/4$ and
$\zeta(3,1,1)=2\zeta(5)-\zeta(2)\zeta(3)$.

\subsection{An all-depth cyclotomic projection}

Use the convention
\[
 \operatorname{Li}_{p,\{1\}^r}(z_0,z_1,\ldots,z_r)
 =\sum_{m_0>m_1>\cdots>m_r\geq1}
   \frac{z_0^{m_0}z_1^{m_1}\cdots z_r^{m_r}}
        {m_0^p m_1\cdots m_r}.
\]
For $p\geq2$ this series is absolutely convergent when all the
$z_j$ are roots of unity.

\begin{corollary}[Explicit reduction of a cyclotomic average]
\label{harm:cor:cyclotomic}
Let $q\geq1$, $1\leq h\leq q$, $\omega=e^{2\pi i/q}$, and
$a=h/q-1$.  For every $p\geq2$ and $r\geq0$,
\begin{equation}\label{harm:eq:projection}
 \sum_{j_0,\ldots,j_r=0}^{q-1}\omega^{-hj_0}
 \operatorname{Li}_{p,\{1\}^r}
       (\omega^{j_0},\ldots,\omega^{j_r})
 =q^{1-p}\frac{(-1)^p}{(p-1)!}
            h_{r+1}^{(p-1)}(h/q-1).
\end{equation}
Every value on the right reduces to depth-one polylogarithms at
$q$th roots of unity and logarithms.  This is a reduction of the
specified average; it makes no assertion that its individual colored
summands all reduce to depth one.
\end{corollary}
\begin{proof}
In \eqref{harm:eq:E}, set $m_0=q(n-1)+h$ and $m_i=qk_i$ for the
indices $n-1\geq k_1>\cdots>k_r\geq1$ in $e_r(n-1)$.  This gives
\[
 E_r(p;h/q-1)=q^{p+r}
 \sum_{\substack{m_0>\cdots>m_r\geq1\\
                  m_0\equiv h\ (q),\ m_i\equiv0\ (q),\ i\geq1}}
 \frac1{m_0^p m_1\cdots m_r}.
\]
Insert the finite Fourier filters for the $r+1$ congruences and
apply \eqref{harm:eq:positive}.  To make the depth-one reduction
explicit, one has, for $k\geq2$,
\begin{align}
 \zeta(k,h/q)
 &=q^{k-1}\sum_{j=0}^{q-1}\omega^{-hj}
                       \operatorname{Li}_k(\omega^j),
                       \label{harm:eq:fourierZ}\\
 H_{h/q-1}^{(1)}
 &=-\log q+\sum_{j=1}^{q-1}\omega^{-hj}\log(1-\omega^j).
                       \label{harm:eq:fourierH}
\end{align}
The first identity follows by grouping a Dirichlet series into residue
classes.  Expand the same identity at spectral argument $1$ to obtain
the second, using $\operatorname{Li}_1(z)=-\log(1-z)$ with principal
logarithms.  Complex conjugate terms have conjugate logarithms, so
\eqref{harm:eq:fourierH} is real.  Formula \eqref{harm:eq:Hcomplex}
and the recurrence \eqref{harm:eq:hrec} finish the reduction.
\end{proof}

For a concrete nonlinear family, let $L=\log2$.  Specializing to
$a=-1/2$ gives
\begin{align}
 \sum_{n\geq0}\frac{H_n}{(2n+1)^2}
 &=\frac{7\zeta(3)-\pi^2L}{4},\label{harm:eq:odd1}\\
 \sum_{n\geq0}\frac{H_n^2-H_n^{(2)}}{(2n+1)^2}
 &=\frac{\pi^2L^2}{2}-7L\zeta(3)+\frac{\pi^4}{24},
                       \label{harm:eq:odd2}\\
 \sum_{n\geq0}\frac{H_n^3-3H_nH_n^{(2)}+2H_n^{(3)}}{(2n+1)^2}
 &=\frac32\bigl[31\zeta(5)-13\zeta(2)\zeta(3)-15L\zeta(4)
             \notag\\
 &\hspace{1.2cm}+14L^2\zeta(3)-4L^3\zeta(2)\bigr].
                       \label{harm:eq:odd3}
\end{align}
All three are instances of one identity, valid as a germ at $u=0$:
\begin{equation}\label{harm:eq:oddall}
 \sum_{r\geq0}u^r\sum_{n\geq0}\frac{e_r(n)}{(2n+1)^2}
 =\frac{\sqrt\pi\,\Gamma(1-u)}{4u\,\Gamma(1/2-u)}
       \bigl[\psi(1/2)-\psi(1/2-u)\bigr].
\end{equation}
This yields arbitrary degree in the harmonic numbers, with no
integer-relation search.  Assigning weight $k$ to $H_a^{(k)}$ and
$\zeta(k,A)$ makes \eqref{harm:eq:positive} homogeneous of weight
$p+r$; the rational Fourier reduction preserves this weight.

\subsection{Meromorphic continuation with a uniform remainder}

The same family gives exact spectral data beyond the half-plane where
its Dirichlet series converges.  Define polynomials $b_k(u,a)$ by
\begin{equation}\label{harm:eq:bdef}
 \phi(t,u,a):=e^{-At}
       \left(\frac{t}{1-e^{-t}}\right)^{1+u}
       =\sum_{k\geq0}b_k(u,a)t^k.
\end{equation}
The branch at $t=0$ is the analytic one taking value $1$.  Equivalently,
in the standard generalized Bernoulli-polynomial notation,
\[
 b_k(u,a)=\frac{B_k^{(1+u)}(u-a)}{k!}.
\]
It is enough to use the defining series; no properties of generalized
Bernoulli polynomials are assumed.  In particular,
\begin{align}
 b_0&=1,\\
 b_1&=\frac{u-2a-1}{2},\label{harm:eq:b1}\\
 b_2&=\frac{12a^2-12au+12a+3u^2-7u+2}{24},\label{harm:eq:b2}\\
 b_3&=\frac{(u-2a-1)(4a^2-4au+4a+u^2-3u)}{48}.
                       \label{harm:eq:b3}
\end{align}
Each $b_k$ is a rational polynomial in $a,u$, of degree at most $k$
in each variable.

\begin{lemma}[Uniform Mellin subtraction]\label{harm:lem:subtract}
Fix a compact subset of $\Re a>-1$ and a sufficiently small disc of
$u$-values around zero.  For every integer $K\geq1$,
\begin{equation}\label{harm:eq:subtraction}
 \Gamma(s)F(s,u;a)
   =\sum_{k=0}^{K-1}\frac{b_k(u,a)}{s-1-u+k}
       +J_K(s,u;a),
\end{equation}
where $J_K$ is jointly holomorphic on
$\Re s>1+\Re u-K$.  Therefore $F$ extends meromorphically to all
$s$, locally in $u$, and $E_r(s;a)=[u^r]F(s,u;a)$ has possible
poles only at $s=1,0,-1,-2,\ldots$.
\end{lemma}
\begin{proof}
Split \eqref{harm:eq:mellinF} at $t=1$, and define
\begin{align*}
 J_K(s,u;a)={}&\int_0^1 t^{s-2-u}
       \left(\phi(t,u,a)-\sum_{k=0}^{K-1}b_k(u,a)t^k\right)dt\\
 &+\int_1^{\infty}t^{s-1}e^{-At}
                  (1-e^{-t})^{-1-u}\,dt.
\end{align*}
The parenthesized remainder is $O(t^K)$ uniformly on the fixed
parameter sets.  Every derivative in $s$ or $u$ introduces only a
fixed power of $\log t$ near zero, so it remains integrable on compact
subsets of $\Re s>1+\Re u-K$.  The second integral and all its
parameter derivatives decay exponentially at infinity because
$\Re A$ stays positive.  Differentiation under both integrals proves
joint holomorphy.  Integrating each subtracted monomial gives
\eqref{harm:eq:subtraction}.  Increasing $K$ continues the same
function, since all formulas agree on the original convergence
domain.  Finally, $1/\Gamma(s)$ is entire; coefficient extraction
in $u$ can turn the simple moving denominators into higher poles
at the stated integers, and introduces no others.
\end{proof}

\subsection{The complete principal part and finite part at \texorpdfstring{$s=1$}{s=1}}

Write
\begin{equation}\label{harm:eq:c}
 C(u)=\frac1{\Gamma(1+u)}=\sum_{j\geq0}c_ju^j
 =\exp\left(\gamma u+
       \sum_{k\geq2}\frac{(-1)^{k+1}\zeta(k)}k u^k\right).
\end{equation}
Thus $c_0=1$, $c_1=\gamma$,
$c_2=(\gamma^2-\zeta(2))/2$, and
$c_3=(\gamma^3-3\gamma\zeta(2)+2\zeta(3))/6$.
The finite part $\operatorname{FP}_{s=s_0}$ always means the
coefficient of $(s-s_0)^0$ in the Laurent expansion of the specified
meromorphic function.

\begin{theorem}[All depths at the principal pole]\label{harm:thm:one}
For $r\geq0$ and $\Re a>-1$, as $\varepsilon\to0$,
\begin{equation}\label{harm:eq:one}
 \boxed{\quad
 E_r(1+\varepsilon;a)
 =\sum_{j=0}^{r}\frac{c_j}{\varepsilon^{r+1-j}}
       +c_{r+1}-h_{r+1}(a)+O(\varepsilon).
 \quad}
\end{equation}
In particular, all principal coefficients are independent of the
shift, while the complete shift dependence of the finite part is
the single complete harmonic polynomial $-h_{r+1}(a)$.
\end{theorem}
\begin{proof}
First take $a=0$.  With $K=1$, the singular term in
\eqref{harm:eq:subtraction} near $(s,u)=(1,0)$ is $1/(s-1-u)$.
Consequently
\[
 F(s,u;0)=\frac{1}{\Gamma(1+u)(s-1-u)}+Q(s,u),
\]
with $Q$ jointly holomorphic there: replacing $1/\Gamma(s)$ by
$1/\Gamma(1+u)$ in the singular numerator changes the expression
by a holomorphic divided difference.  For $s=1$ and small
$\Re u<0$, the beta integral gives $F(1,u;0)=-1/u$, and hence
\[
 Q(1,u)=\frac{C(u)-1}{u}.
\]
The right side extends holomorphically to $u=0$.  Extracting
$[u^r]$ from $C(u)/(\varepsilon-u)$ gives exactly
$\sum_{j=0}^{r}c_j\varepsilon^{-r-1+j}$, and extracting it from
$Q(1,u)$ gives $c_{r+1}$.

For arbitrary $a$, the difference
\[
 E_r(s;a)-E_r(s;0)
   =\sum_{n\geq1}e_r(n-1)
               \bigl((n+a)^{-s}-n^{-s}\bigr)
\]
is normally convergent for $\Re s>0$, locally uniformly in $a$.
At $s=1$ it equals $-h_{r+1}(a)$ by
\eqref{harm:eq:primitive}.  This leaves the principal part unchanged
and gives the stated finite part.
\end{proof}

For example,
\begin{align}
 \operatorname{FP}_{s=1}E_0(s;a)&=-\psi(A),\\
 \operatorname{FP}_{s=1}E_1(s;a)
 &=\frac{\gamma^2-\zeta(2)-(H_a^{(1)})^2-H_a^{(2)}}2,\\
 \operatorname{FP}_{s=1}E_2(s;a)
 &=\frac{\gamma^3-3\gamma\zeta(2)+2\zeta(3)}6-h_3(a).
\end{align}
The first line is the usual zeroth generalized Stieltjes constant.
For higher depths, the corresponding zeroth constants are just as
explicit.  This does not imply such a reduction for all higher
Laurent coefficients.

\subsection{Every nonpositive integer: a finite exact coefficient rule}

\begin{theorem}[Principal and finite Laurent data at nonpositive integers]
\label{harm:thm:negative}
Let $m,r\geq0$, and define
\begin{equation}\label{harm:eq:A}
 \mathcal A_m(u,a)=\frac{b_{m+1}(u,a)}{\Gamma(u-m)}
                =\sum_{d\geq1}A_{m,d}(a)u^d.
\end{equation}
Then
\begin{equation}\label{harm:eq:negative}
 \boxed{\quad
 E_r(-m+\varepsilon;a)
  =\sum_{h=0}^{r}A_{m,r+1-h}(a)\varepsilon^{-h}
        +O(\varepsilon).
 \quad}
\end{equation}
In particular,
\begin{equation}\label{harm:eq:FPnegative}
 \operatorname{FP}_{s=-m}E_r(s;a)
       =[u^{r+1}]\frac{b_{m+1}(u,a)}{\Gamma(u-m)}.
\end{equation}
All coefficients displayed in \eqref{harm:eq:negative} belong to
$\mathbb Q[a,\gamma,\zeta(2),\ldots,\zeta(r)]$, with the obvious
interpretation for $r=0,1$.  The highest possible pole coefficient is
\begin{equation}\label{harm:eq:leadingpole}
 A_{m,1}(a)=\zeta(-m,A).
\end{equation}
\end{theorem}
\begin{proof}
Choose $K>m+1$ in Lemma~\ref{harm:lem:subtract}.  Near
$(s,u)=(-m,0)$, precisely one of its denominators can vanish:
the one with $k=m+1$.  Thus, putting $s=-m+\varepsilon$,
\[
 F(-m+\varepsilon,u;a)
  =\frac{b_{m+1}(u,a)}{\Gamma(-m+\varepsilon)(\varepsilon-u)}
       +\frac{Q_m(\varepsilon,u;a)}{\Gamma(-m+\varepsilon)},
\]
where $Q_m$ is jointly holomorphic.  Since $1/\Gamma(-m+\varepsilon)$
has a simple zero at $\varepsilon=0$, the second term contributes
$O(\varepsilon)$ after extracting any fixed coefficient in $u$.
Write
\[
 b_{m+1}(u,a)=\sum_{j\geq0}b_{m+1,j}(a)u^j,
 \qquad \frac1{\Gamma(-m+\varepsilon)}
              =\sum_{\ell\geq1}g_{m,\ell}\varepsilon^\ell.
\]
Coefficient extraction from the singular term gives
\[
 \left(\sum_{\ell\geq1}g_{m,\ell}\varepsilon^\ell\right)
       \sum_{j=0}^r b_{m+1,j}(a)\varepsilon^{-r+j-1}.
\]
Its coefficient at $\varepsilon^{-h}$ is
$\sum_{j=0}^{r-h}b_{m+1,j}(a)g_{m,r+1-h-j}$, which is exactly
$[u^{r+1-h}]\mathcal A_m(u,a)$.  This proves
\eqref{harm:eq:negative}.

The reciprocal-gamma factor has the explicit harmonic expansion
\begin{equation}\label{harm:eq:recipgamma}
 \frac1{\Gamma(u-m)}=(-1)^m m!\,u
 \exp\left((\gamma-H_m)u+
   \sum_{k\geq2}\frac{(-1)^{k+1}\zeta(k)-H_m^{(k)}}k u^k\right).
\end{equation}
Indeed, multiply $1/\Gamma(1+u)$ by
$u\prod_{j=1}^m(u-j)$ and expand its logarithm after removing
$(-1)^m m!u$.  The finite coefficient rule now proves the claimed
polynomial ring.  Finally,
$b_{m+1}(0,a)=B_{m+1}(-a)/(m+1)!$ and
$g_{m,1}=(-1)^m m!$.  Bernoulli reflection and
$\zeta(-m,A)=-B_{m+1}(A)/(m+1)$ give
\eqref{harm:eq:leadingpole}.
\end{proof}

For a generic shift, the pole at $s=-m$ has order $r$ when $r\geq1$.
At a zero of $\zeta(-m,A)$ its order drops.  Thus the familiar
absence of negative even poles for the depth-one unshifted harmonic
zeta function is part of an all-depth cancellation rule: at $a=0$
and positive even $m$, the pole order is at most $r-1$.  The formula
determines whether further cancellations occur.

No divergent series is assigned an ordinary sum here.  The values in
\eqref{harm:eq:FPnegative} are finite Laurent coefficients of a
specified meromorphic continuation.  Nor can they be computed by
first setting $s=-m$ in $F$ at a nonzero $u$ and then differentiating:
the moving divisor $\varepsilon-u=0$ makes those operations
incompatible at their intersection.

\subsection{Small exact tables and differentiation of finite parts}

At $s=0$, \eqref{harm:eq:b1} makes the whole finite-part family
particularly simple.  With $c_{-1}=0$,
\begin{equation}\label{harm:eq:zeroall}
 \operatorname{FP}_{s=0}E_r(s;a)
       =\frac12c_{r-1}-\left(a+\frac12\right)c_r.
\end{equation}
More generally every finite part at $s=-m$ is a polynomial in $a$
of degree at most $m+1$.  The first unshifted examples are shown in
Table~\ref{harm:tab:FP}; only finite parts, not complete Laurent
series, are tabulated.

\begin{table}[htbp]
\centering
\small
\renewcommand{\arraystretch}{1.6}
\begin{tabular}{c|ccc}
 & $s=0$ & $s=-1$ & $s=-2$\\\hline
 $r=0$ & $-\tfrac12$ & $-\tfrac1{12}$ & $0$\\
 $r=1$ & $\tfrac{1-\gamma}{2}$ & $\tfrac{9-2\gamma}{24}$ & $\tfrac18$\\
 $r=2$ & $\tfrac{-\gamma^2+2\gamma+\zeta(2)}4$
       & $\tfrac{-\gamma^2+9\gamma+\zeta(2)-10}{24}$
       & $\tfrac{6\gamma-17}{48}$
\end{tabular}
\caption{Exact finite parts of $E_r(s;0)$ from the single
coefficient rule \eqref{harm:eq:FPnegative}.}
\label{harm:tab:FP}
\end{table}

For example, the complete displayed portion at a point where the
highest possible pole cancels is
\begin{align}
 E_3(-2+\varepsilon;0)
  ={}&\frac1{8\varepsilon^2}
       +\frac{6\gamma-17}{48\varepsilon}
       +\frac{3\gamma^2-17\gamma-3\zeta(2)+17}{48}
       +O(\varepsilon).\label{harm:eq:cancelledexample}
\end{align}

There is a necessary residue correction when finite parts are
differentiated.  Denote
\[
 C_{m,r}(a)=\operatorname{FP}_{s=-m}E_r(s;a)\quad(m\geq0),
 \qquad C_{-1,r}(a)=\operatorname{FP}_{s=1}E_r(s;a),
\]
and let $R_{m,r}(a)$ denote the residue at $s=-m$, also for $m=-1$.

\begin{corollary}[Differentiation with its residue term]
\label{harm:cor:residue}
For $m\geq0$,
\begin{equation}\label{harm:eq:residuecorrection}
 C_{m,r}'(a)=mC_{m-1,r}(a)-R_{m-1,r}(a).
\end{equation}
In particular $C_{0,r}'(a)=-c_r$.
\end{corollary}
\begin{proof}
Normal convergence first gives
$\partial_a E_r(s;a)=-sE_r(s+1;a)$ for $\Re s>1$;
Lemma~\ref{harm:lem:subtract} continues the identity meromorphically.
Put $s=-m+\varepsilon$.  The constant term of
$(m-\varepsilon)E_r(1-m+\varepsilon;a)$ is
$mC_{m-1,r}-R_{m-1,r}$, even if the latter function has higher
order poles.  Taking the constant coefficient commutes with
$a$-differentiation by joint meromorphy on compact shift domains.
At $m=0$, Theorem~\ref{harm:thm:one} gives $R_{-1,r}=c_r$,
in agreement with \eqref{harm:eq:zeroall}.
\end{proof}

There is also a direct coefficient-level verification.  Since
$\partial_a b_{m+1}(u,a)=-b_m(u,a)$ and
$\Gamma(u-m+1)=(u-m)\Gamma(u-m)$,
\[
 \partial_a\mathcal A_m(u,a)=(m-u)\mathcal A_{m-1}(u,a),
 \qquad m\geq1.
\]
Taking the coefficient of $u^{r+1}$ gives exactly
\eqref{harm:eq:residuecorrection}; the case $m=0$ is
\eqref{harm:eq:zeroall}.

Dropping the residue in \eqref{harm:eq:residuecorrection} would,
already at $m=0$, incorrectly predict a constant function of $a$.
This is a correction to a tempting manipulation, rather than an
assertion that the inspected manuscript makes that error.

\subsection{Higher Stieltjes coefficients and the precise scope of reduction}

Define the higher shifted height-one constants by
\begin{equation}\label{harm:eq:Stieltjes}
 E_r(1+\varepsilon;a)
 =\sum_{j=0}^{r}c_j\varepsilon^{-r-1+j}
     +\sum_{\ell\geq0}\frac{(-1)^\ell}{\ell!}
                   \gamma_{r,\ell}(a)\varepsilon^\ell.
\end{equation}
For $r=0$, these are the usual generalized Stieltjes constants
$\gamma_{0,\ell}(a)=\gamma_\ell(A)$.
Theorem~\ref{harm:thm:one} evaluates
$\gamma_{r,0}(a)=c_{r+1}-h_{r+1}(a)$.
For real $a,b>-1$ and every $\ell\geq0$, normal convergence of the
shifted difference gives the exact, absolutely convergent identity
\begin{equation}\label{harm:eq:shiftStieltjes}
 \gamma_{r,\ell}(a)-\gamma_{r,\ell}(b)
  =\sum_{n\geq1}e_r(n-1)
       \left(\frac{\log^\ell(n+a)}{n+a}
                  -\frac{\log^\ell(n+b)}{n+b}\right).
\end{equation}
The summand is $O(n^{-2}(1+\log^{r+\ell}n))$.  At $\ell=0$
this reduces to $h_{r+1}(b)-h_{r+1}(a)$.  At higher $\ell$ it
specifies the shift dependence without assuming a finite reduction
of the individual constants.

One also obtains the all-depth extension of the manuscript's
Stieltjes differentiation ladder.  Let
$e_i(k)=[v^i]\prod_{j=1}^k(1+v/j)$ and use $\partial_s$ for
spectral differentiation.  For $k\geq1$,
\begin{equation}\label{harm:eq:Stieltjesderivative}
 \frac{d^k}{da^k}\gamma_{r,\ell}(a)
 =(-1)^{\ell+k}k!\ell!
      \sum_{i+j=\ell}e_i(k)
             \frac{\partial_s^j E_r(k+1;a)}{j!}.
\end{equation}
Indeed, $\partial_a^kE_r(s;a)=(-1)^k(s)_kE_r(s+k;a)$;
the principal part at $s=1$ is independent of $a$, and coefficient
comparison at $s=1$ proves the formula.  When $\ell=0$ the right
side is fully evaluated by Theorem~\ref{harm:thm:beta}.  For
$\ell>0$ it involves genuine spectral derivatives.

For completeness these derivatives have exact logarithmic Mellin
representations.  Differentiating \eqref{harm:eq:mellinE} yields
\begin{align}
 &\frac1{r!}\int_0^1
   \frac{x^a(-\log x)^{s-1}[-\log(1-x)]^r}{1-x}
                     \log^\ell(-\log x)\,dx\notag\\
 &\hspace{1cm}
 =\sum_{j=0}^{\ell}\binom{\ell}{j}
      \Gamma^{(\ell-j)}(s)\,\partial_s^jE_r(s;a),
      \qquad \Re s>1.\label{harm:eq:logspectral}
\end{align}
The endpoint estimates in Proposition~\ref{harm:prop:mellin} remain
integrable after any fixed number of logarithmic factors, justifying
all differentiations.  Gamma derivatives here can equivalently be
written using complete Bell polynomials in polygamma functions.

The exact polynomial conclusion of
Theorem~\ref{harm:thm:negative} stops at the constant Laurent
coefficient.  The term $Q_m/\Gamma(s)$ in its proof starts at
order $s+m$ and can contribute new constants from that point
onward.  At depth zero these already include $\zeta'(-m,A)$,
and at $m=0$ Lerch's identity gives the log-gamma function.
There is therefore no basis for discarding the analytic remainder
when computing positive spectral derivatives.

The next concrete questions are to identify which linear combinations
of these higher coefficients eliminate the analytic remainder, to
extend the cyclotomic projection to harmonic patterns other than
$\{1\}^r$, and to compare the outer-shift constants
\eqref{harm:eq:Stieltjes} systematically with the literature's
simultaneously shifted harmonic Hurwitz constants.  Those extensions
are left open here.

% END input


% BEGIN sections/04b_harmonic_symmetries.tex
\subsection{Reflection and antiderivatives of the harmonic finite parts}
\label{harm:subsec:symmetries}

The shifted coefficient theorem also gives exact transformation laws
that are not visible in an unshifted table. Put
\[
 C_{m,r}(a)=\operatorname{FP}_{s=-m}E_r(s;a)
          =[u^{r+1}]\mathcal A_m(u,a),\qquad m,r\geq0,
\]
and set \(C_{m,r}=0\) for \(r<0\).
Each \(C_{m,r}\) is a polynomial in \(a\), of degree at most \(m+1\).
Thus the identities below hold as polynomial identities for all complex
\(a\); they refer directly to the original shifted Dirichlet family
where \(\Re a>-1\).

\begin{theorem}[Reflection mixes harmonic depths]
\label{harm:thm:reflection}
For all \(m,r\geq0\),
\begin{equation}\label{harm:eq:FPreflection}
 C_{m,r}(-1-a)
 =(-1)^{m+1}\sum_{j=0}^{r}
          \frac{1}{j!}\frac{d^j}{da^j}C_{m,r-j}(a).
\end{equation}
Terms with \(j>m+1\) vanish. On the strip \(-1<\Re a<0\),
both sides are finite parts of the original shifted families.
Elsewhere the left side means the polynomial continuation of its
finite-part coefficient.
\end{theorem}

\begin{proof}
The defining local kernel can be written
\[
 \phi(t,u,a)=
 e^{((u-1)/2-a)t}
 \left(\frac{t}{2\sinh(t/2)}\right)^{1+u}.
\]
All powers here mean their Taylor germs at \(t=0\), where the
parenthesized factor equals one. Its second factor is even in \(t\).
Consequently
\[
 \phi(t,u,-1-a)=\phi(-t,u,a+u),\qquad
 b_{m+1}(u,-1-a)=(-1)^{m+1}b_{m+1}(u,a+u).
\]
The Gamma denominator is independent of \(a\), so the same identity
holds for \(\mathcal A_m\). Expanding its polynomial argument
translation in powers of \(u\) and taking \([u^{r+1}]\) proves
\eqref{harm:eq:FPreflection}. There is no constant coefficient of
\(\mathcal A_m\), which is why the sum ends at \(j=r\).
\end{proof}

For \(r=0\), this is the familiar Bernoulli reflection
\(\zeta(-m,-a)=(-1)^{m+1}\zeta(-m,a+1)\), interpreted as a polynomial
identity. At positive depth the derivative terms are essential.
For example,
\[
 C_{1,1}(-1-a)-C_{1,1}(a)
 =C_{1,0}'(a)=-a-\frac12.
\]
Thus an ordinary Bernoulli parity rule at fixed depth would be false.

\begin{theorem}[Complete finite-part antiderivative ladder]
\label{harm:thm:FPprimitives}
Let \(q\geq1\), and define \(h_j^{[m,q]}\) by
\begin{equation}\label{harm:eq:primitiveh}
 \prod_{\nu=1}^{q}
       \left(1-\frac{u}{m+\nu}\right)^{-1}
       =\sum_{j\geq0}h_j^{[m,q]}u^j.
\end{equation}
Then the polynomial
\begin{equation}\label{harm:eq:FPprimitive}
 F_{m,r}^{(q)}(a)=
 \frac{1}{(m+1)_q}
 \sum_{j=0}^{r}h_j^{[m,q]}C_{m+q,r-j}(a)
\end{equation}
satisfies
\[
 \frac{d^q}{da^q}F_{m,r}^{(q)}(a)=C_{m,r}(a).
\]
Here \((m+1)_q=(m+1)(m+2)\cdots(m+q)\).
Every \(q\)-fold antiderivative differs from this one by a polynomial
of degree less than \(q\).
\end{theorem}

\begin{proof}
Differentiating the local kernel gives
\(\partial_a b_k=-b_{k-1}\), hence
\[
 \partial_a\mathcal A_m=(m-u)\mathcal A_{m-1}.
\]
Iteration, in the direction of increasing \(m\), gives
\[
 \partial_a^q\mathcal A_{m+q}
       =\prod_{\nu=1}^{q}(m+\nu-u)\,\mathcal A_m.
\]
Divide by this polynomial, which is invertible as a germ at \(u=0\),
and extract \([u^{r+1}]\). Equation
\eqref{harm:eq:primitiveh} gives exactly
\eqref{harm:eq:FPprimitive}.
\end{proof}

The coefficients in this ladder are again generalized-harmonic
Bell coefficients:
\[
 h_j^{[m,q]}=[u^j]\exp\left(
   \sum_{k\geq1}\frac{H_{m+q}^{(k)}-H_m^{(k)}}{k}u^k
                              \right).
\]
In particular,
\begin{equation}\label{harm:eq:firstFPprimitive}
 \int C_{m,r}(a)\,da
 =\sum_{j=0}^{r}\frac{C_{m+1,r-j}(a)}{(m+1)^{j+1}}
   +\text{constant}.
\end{equation}
The lower-depth terms are exactly the residue correction propagated
through integration. To normalize a repeated integral at \(a_0\), use
\[
 F_{m,r}^{(q)}(a)
 -\sum_{j=0}^{q-1}
       \frac{(a-a_0)^j}{j!}
       \left.\frac{d^j}{da^j}F_{m,r}^{(q)}(a)\right|_{a=a_0}.
\]
This polynomial and its first \(q-1\) derivatives vanish at \(a_0\);
it is the \(q\)-fold integral with those initial conditions.

\subsection{Rational multiplication with explicit depth corrections}
\label{harm:subsec:multiplication}

An ordinary Hurwitz multiplication identity does not survive unchanged
when the outer denominator is shifted but the inner harmonic
denominators are fixed. Its finite-part replacement is nevertheless a
finite exact formula.

For an integer \(q\geq2\), let
\begin{equation}\label{harm:eq:dq}
 \begin{split}
 \Lambda_q(t)&=\log\left(
       \frac{1-e^{-t}}{q(1-e^{-t/q})}\right),\\
 d_n^{(q)}(u)&=[t^n]\exp\bigl(-u\Lambda_q(t)\bigr).
 \end{split}
\end{equation}
The logarithm is its analytic germ with value zero at \(t=0\).
Each \(d_n^{(q)}\) is a rational polynomial in \(u\) of degree at most
\(n\), and \(d_n^{(q)}(0)=0\) for \(n>0\). Set
\[
 \mathcal A_{-1}(u,a)=\frac{1}{\Gamma(1+u)}
                    =\sum_{j\geq0}c_j u^j.
\]
This auxiliary term is independent of \(a\); it is not the generating
function of the finite parts at \(s=1\).

\begin{theorem}[Multiplication law for the singular coefficient germ]
\label{harm:thm:multiplication}
For \(m\geq0\), put \(a_j=(a+1+j)/q-1\). Then
\begin{equation}\label{harm:eq:multiplication}
 \sum_{j=0}^{q-1}\mathcal A_m(u,a_j)
 =\sum_{n=0}^{m+1}
       q^{\,n-m}d_n^{(q)}(u)(u-m)_n
                     \mathcal A_{m-n}(u,a).
\end{equation}
All factors are analytic germs at \(u=0\). Taking \([u^{r+1}]\)
gives a finite multiplication formula for \(C_{m,r}\); taking lower
powers gives the corresponding principal Laurent coefficients.
\end{theorem}

\begin{proof}
Sum the defining local kernels at the indicated shifts:
\[
 \begin{split}
 \sum_{j=0}^{q-1}\phi(t,u,a_j)
 &=e^{-(a+1)t/q}
       \frac{1-e^{-t}}{1-e^{-t/q}}
       \left(\frac{t}{1-e^{-t}}\right)^{1+u}\\
 &=q\,\phi(t/q,u,a)\exp\bigl(-u\Lambda_q(t)\bigr).
 \end{split}
\]
Taking \([t^{m+1}]\) and dividing by \(\Gamma(u-m)\) gives
\[
 \sum_{j=0}^{q-1}\mathcal A_m(u,a_j)
 =\sum_{n=0}^{m+1}q^{\,n-m}d_n^{(q)}(u)
       \frac{b_{m+1-n}(u,a)}{\Gamma(u-m)}.
\]
The Gamma recurrence identifies the final quotient as
\((u-m)_n\mathcal A_{m-n}\), including \(n=m+1\), since
\(b_0=1\). This proves the identity without interchanging a
singular specialization with an infinite sum.
\end{proof}

At depth zero the correction terms vanish, and the formula reduces to
\[
 \sum_{j=0}^{q-1}
       \zeta\left(-m,\frac{a+1+j}{q}\right)
       =q^{-m}\zeta(-m,a+1).
\]
For positive depth, the finite correction terms measure the failure of
that unmodified rule.

The first instance is especially simple:
\begin{equation}\label{harm:eq:multiplicationzero}
 \sum_{j=0}^{q-1}C_{0,r}(a_j)
       =C_{0,r}(a)+\frac{q-1}{2}c_{r-1},
 \qquad c_j=0\quad(j<0).
\end{equation}
Indeed \(d_1^{(q)}(u)=u(q-1)/(2q)\). For \(q=2,m=1\),
\[
 d_1^{(2)}(u)=\frac u4,\qquad
 d_2^{(2)}(u)=\frac{u^2-u}{32},
\]
so
\begin{equation}\label{harm:eq:duplicationgerm}
 \mathcal A_1\left(u,\frac{a-1}{2}\right)
 +\mathcal A_1\left(u,\frac a2\right)
 =\frac12\mathcal A_1(u,a)
   +\frac{u(u-1)}4\mathcal A_0(u,a)
   +\frac{u^2(u-1)^2}{16\Gamma(1+u)}.
\end{equation}
For example, its depth-one constant term gives
\[
 C_{1,1}\left(\frac{a-1}{2}\right)
 +C_{1,1}\left(\frac a2\right)
 =\frac12C_{1,1}(a)+\frac a4+\frac3{16}.
\]

\subsection{Half-shift Laurent identities}

The coefficient rule and its duplication law give concise values at the
half shift. With \(a=-1/2\),
\[
 b_2(u,-1/2)=\frac{3u^2-u-1}{24},\qquad
 b_3(u,-1/2)=\frac{u(u^2-u-1)}{48}.
\]
Combining these with \(1/\Gamma(u-m)\) yields
\begin{align}
 E_1(-1+\varepsilon;-1/2)
  &=\frac{1}{24\varepsilon}+\frac{\gamma}{24}
       +O(\varepsilon),\label{harm:eq:half11}\\
 E_2(-1+\varepsilon;-1/2)
  &=\frac{1}{24\varepsilon^2}
      +\frac{\gamma}{24\varepsilon}
      +\frac{\gamma^2-\zeta(2)-8}{48}
       +O(\varepsilon),\label{harm:eq:half21}\\
 E_1(-2+\varepsilon;-1/2)
  &=-\frac1{24}+O(\varepsilon),\label{harm:eq:half12}\\
 E_2(-2+\varepsilon;-1/2)
  &=-\frac{1}{24\varepsilon}+\frac{1-2\gamma}{48}
       +O(\varepsilon).\label{harm:eq:half22}
\end{align}
For example, \(E_1(s;-1/2)=\sum_{n\geq1}H_{n-1}(n-1/2)^{-s}\)
in its convergent half-plane. The displayed values mean its
meromorphic continuation, not convergent sums at negative \(s\).
The drop in pole order at \(-2\) follows from the vanishing
\(\zeta(-2,1/2)=0\); its next coefficient does not vanish.
These formulas illustrate exact cancellation across depth, shift, and
spectral order.


% END input


% BEGIN sections/05_polylogarithms.tex
\section{An explicit polylogarithm realization of the convolution calculus}
\label{poly:sec}

The circle Fourier description gives a direct bridge to the
polylogarithm's order derivatives at nonpositive integers. The underlying
Fourier summation is classical \cite{DLMFPolylog,DLMFHurwitz}.
The full mixed kernel and the symmetric first-Stieltjes specialization
also occur in the archived continuations
\cite{RepoConvolutionAlgebra,RepoConvolutionCalculus}.
This section consolidates their order-derivative form while preserving
the endpoint distributions.

\subsection{A two-parameter kernel}

For \(\alpha,\beta\in\bbC\), define the zero-mean distribution
\(\mathcal N_{\alpha,\beta}\) by
\begin{equation}\label{poly:N-fourier}
 \widehat{\mathcal N_{\alpha,\beta}}(n)=
 (2\pi\numi n)^\alpha(-2\pi\numi n)^\beta\quad(n\ne0),
 \qquad \widehat{\mathcal N_{\alpha,\beta}}(0)=0.
\end{equation}
The two powers use the logarithms specified in Section~\ref{scope:sec}.
On every compact set of parameters their coefficients and any fixed
number of parameter derivatives have a common polynomial growth
bound. This proves that \(\mathcal N\) is entire as a
distribution-valued function.

\begin{theorem}[Polylogarithm kernel for every mixed order]
\label{poly:master}
For \(0<h<1\), put \(z=\rme^{2\pi\numi h}\). Then
\begin{align}\label{poly:N}
 \mathcal N_{\alpha,\beta}(h)
 =(2\pi)^{\alpha+\beta}\bigl[
 &\rme^{\pi\numi(\alpha-\beta)/2}
                 \rLi_{-\alpha-\beta}(z)\notag\\
 &+\rme^{-\pi\numi(\alpha-\beta)/2}
                 \rLi_{-\alpha-\beta}(z^{-1})\bigr].
\end{align}
The formula represents the restriction to the open interval of the
entire distribution family, not an assertion of ordinary convergence
of its Fourier series at every order.

Let \(P_k\) be the spectral finite-part polygamma distribution, so that
\[
 \widehat P_k(n)=(2\pi\numi n)^k
  \bigl[\gamma+\log(2\pi\numi n)-H_k\bigr]\quad(n\ne0),
 \qquad H_0=0.
\]
For all \(k,l\geq0\), one has the distribution identity
\begin{equation}\label{poly:mixed}
 P_k*\check P_l=
 \left.
 (\partial_\alpha+\gamma-H_k)
 (\partial_\beta+\gamma-H_l)
       \mathcal N_{\alpha,\beta}
 \right|_{\alpha=k,\ \beta=l}.
\end{equation}
For the derivative-compatible normalization \(Q_k=D^kP_0\), the same
formula holds with \(H_k,H_l\) replaced by zero.
\end{theorem}

\begin{proof}
If \(\Re(\alpha+\beta)<-1\), the Fourier series is absolutely convergent.
The positive and negative indices separately give
\eqref{poly:N}. The same argument applies after parameter
differentiation on any smaller open region. The right side has an
entire continuation in its orders for \(z\ne1\), and is smooth in \(h\)
on compact subintervals of \((0,1)\). To identify its distributional
extension, insert the Abel factor \(\rho^{|n|}\), \(0<\rho<1\),
and take its limit in distributions. The Hurwitz Fourier formula
identifies the same local analytic continuation. The identity
therefore extends to all \(\alpha,\beta\). This argument also
specifies terms supported at \(h=0\), which a pointwise formula alone
could not recover.

Each parameter derivative in \eqref{poly:N-fourier} multiplies the
coefficient by the corresponding logarithm. At \((k,l)\), this gives
\(\widehat P_k(n)\widehat P_l(-n)\), exactly the Fourier coefficient of
\(P_k*\check P_l\). Fourier coefficients determine a distribution on
the circle. The formula for \(Q_k\) follows from its Fourier coefficient
\((2\pi\numi n)^k[\gamma+\log(2\pi\numi n)]\).
\end{proof}

The sine phases must be retained for complex parameters. A real-part
notation would not preserve holomorphic dependence on
\(\alpha,\beta\). At integer orders
\(\mathcal N_{k,l}=(-1)^lD^{k+l}\Delta\).
For \(k+l>0\) this is supported at the endpoint; for \(k=l=0\) its
open-interval value is \(-1\). Its parameter derivatives retain
additional, nonconstant interior terms. Substituting an integer
before differentiating would lose the desired information.

The unreflected kernel is obtained using one spectral parameter:
\(\mathcal M_\alpha=\mathcal N_{\alpha,0}\). Then
\begin{align}\label{poly:unreflected}
 P_k*P_l=\left.
 \bigl[\partial_\alpha^2
  +(2\gamma-H_k-H_l)\partial_\alpha
  +(\gamma-H_k)(\gamma-H_l)\bigr]\mathcal M_\alpha
 \right|_{\alpha=k+l}.
\end{align}
The equality is checked on each nonzero Fourier coefficient.
It also determines all the contact terms.

\subsection{First-Stieltjes identities at polylogarithm order zero}

Write \(c=\gamma+\log(2\pi)\), \(d=\pi/2\), and
\[
 L_j(h)=\left.\partial_s^j
                 \rLi_s(\rme^{2\pi\numi h})\right|_{s=0}.
\]
These derivatives are values of the analytically continued function.
They are not defined by directly summing a growing logarithmic
Fourier series.

\begin{corollary}[Two explicit order-derivative reductions]
\label{poly:gamma1}
For every \(0<h<1\),
\begin{align}
 \gamma_1(h)+\gamma_1(1-h)
 &=2\Re L_2(h)-4c\Re L_1(h)-c^2+\frac{\pi^2}{12},
 \label{poly:symmetric}\\
 \gamma_1(h)
 &=\Re L_2(h)-2c\Re L_1(h)+\pi\Im L_1(h)
      -\frac{c^2}{2}+\frac{\pi^2}{24}
      -\frac{\pi c}{2}\cot(\pi h).
 \label{poly:individual}
\end{align}
\end{corollary}

\begin{proof}
At \(k=l=0\), the operator on the positive-frequency part of
\eqref{poly:N} is
\[
 (c+d\numi-\partial_s)(c-d\numi-\partial_s)
       =(c-\partial_s)^2+d^2 .
\]
For real \(h\) the negative-frequency part is its conjugate.
The reflected convolution identity proved above reads on the open
interval
\[
 \gamma_1(h)+\gamma_1(1-h)-\frac{\pi^2}{3}
 =2\Re\bigl[L_2-2cL_1+(c^2+\pi^2/4)L_0\bigr].
\]
Since \(\Re L_0=-1/2\), this is \eqref{poly:symmetric}.

For the unreflected convolution, the operator is
\((c+d\numi-\partial_s)^2\). Its value is
\(2\gamma_1(h)+\zeta(2)\). Finally,
\[
 L_0(h)=-\frac12+\frac{\numi}{2}\cot(\pi h)
\]
gives \eqref{poly:individual} after expanding the real part.
Every constant comes from an explicit Gamma coefficient or
a contact term.
\end{proof}

These formulas belong to the classical differentiated Hurwitz
functional-equation framework. Their role here is an explicit
application of the finite-part normalization. They also check the
sign of the \(\cot(\pi h)\) contribution and the \(\pi^2\) constants.

At \(h=1/2\), the identity
\(\rLi_s(-1)=(2^{1-s}-1)\zeta(s)\) turns the corollary into an identity
among \(\zeta^{[j]}(0)\), \(\gamma_1\), \(\gamma\), and \(\log2\).
On the Stieltjes side,
\[
 \gamma_1(1/2)=\gamma_1-2\gamma\log2-\log^22,
\]
obtained directly by expanding
\(\zeta(s,1/2)=(2^s-1)\zeta(s)\).
Thus the convergent digamma convolution from the preceding part,
at \(a=1/2\), has the explicit value
\begin{equation}\label{poly:half-convolution}
 2\gamma_1-4\gamma\log2-2\log^22+\frac{\pi^2}{6}.
\end{equation}
This evaluates the fully subtracted integral; the corresponding
unsubtracted product integral diverges.

\subsection{Rational shifts and a finite coordinate algorithm}

The pending report asks for rational-shift coordinate formulas
\cite{RepoIncoming}. The conversion follows from a finite Fourier
transform already present in the manuscript. The distributional
extension now makes it applicable to mixed polygamma kernels at
singular spectral orders.

For \(q\geq2\), \(1\leq p<q\), and \(z=\rme^{2\pi\numi p/q}\),
\begin{equation}\label{poly:rational}
 \rLi_s(z)=q^{-s}\sum_{a=1}^q z^a\zeta(s,a/q).
\end{equation}
At any regular spectral point \(s_0\),
\begin{equation}\label{poly:rational-jets}
 \rLi_{s_0}^{[j]}(z)
 =q^{-s_0}\sum_{a=1}^q z^a
   \sum_{v=0}^j\binom jv(-\log q)^{j-v}
             \zeta^{[v]}(s_0,a/q).
\end{equation}
At \(s=1\) the pole cancels because \(\sum_{a=1}^qz^a=0\);
one must expand the finite combined expression before specializing.
Equations \eqref{poly:N}--\eqref{poly:rational-jets}
are a finite algorithm for every rational-shift polygamma
convolution and every fixed spectral derivative.

A multiplicative-character transform then converts these Hurwitz
coordinates to primitive Dirichlet-\(L\) coordinates. Conductor
logarithms are generated by differentiating the explicit factors.
The manuscript's conductor-lifted trace theorem and pole-corrected
jets already prove that arithmetic conversion \cite{RepoManuscript}.
We use them as an available normal form and do not claim
independence of the resulting evaluated constants.

% END input


% BEGIN sections/06_audit.tex
\section{Source corrections and the scope of the conclusions}
\label{audit:sec:corrections}

The following audit concerns the pinned manuscript, rather than earlier
versions already corrected there.  It distinguishes a branch error, an
unsupported description of a special value, and an editorial clarification
about coefficient fields.  The accompanying patch changes only the first
two items; it does not alter the source files in this research package.

\subsection{A real branch correction for the arctanh integral}

In \path{chapters/03-algebraic.tex}, equation
\texttt{cleo:eq:lintanh} occurs in the discussion of real parameters of
absolute value less than one.  It writes
\[
 \int_0^{\pi/2}\operatorname{arctanh}(\lambda\sin x)\,dx
 =\theta\log\tan\frac\theta2
       +2\operatorname{Cl}_2(\theta)
       -\frac12\operatorname{Cl}_2(2\theta),
 \qquad \theta=\arcsin\lambda.
\]
The displayed logarithm is correct for positive $\lambda$.  For negative
$\lambda$, its argument is negative, and the principal complex logarithm
adds an imaginary term that cannot occur in the real integral.

\begin{proposition}[The real-parameter formula]
\label{audit:prop:arctanh}
For $-1\le\lambda\le1$, put $\theta=\arcsin\lambda$.  Then
\begin{equation}
 \int_0^{\pi/2}\operatorname{arctanh}(\lambda\sin x)\,dx
 =\theta\log\left|\tan\frac\theta2\right|
       +2\operatorname{Cl}_2(\theta)
       -\frac12\operatorname{Cl}_2(2\theta).
 \label{audit:eq:arctanh-corrected}
\end{equation}
At $\lambda=0$ the right side is its continuous value zero.  At
$\lambda=\pm1$ the integrals are convergent improper integrals.
\end{proposition}

\begin{proof}
Let the integral be $F(\lambda)$.  For $|\lambda|<1$, differentiation
under the integral is valid locally uniformly in $\lambda$.  The
substitution $t=\cos x$ gives
\[
 F'(\lambda)=\int_0^1
       \frac{dt}{1-\lambda^2+\lambda^2t^2}
 =\frac{\arcsin\lambda}{\lambda\sqrt{1-\lambda^2}},
\]
where the quotient has value one at zero.  Therefore
$dF(\sin\theta)/d\theta=\theta/\sin\theta$.

The classical Fourier definition of the Clausen function gives
$\operatorname{Cl}_2'(u)=-\log|2\sin(u/2)|$ away from multiples
of $2\pi$.  On either interval $(-\pi/2,0)$ or $(0,\pi/2)$,
differentiating the right side of \eqref{audit:eq:arctanh-corrected}
gives $\theta/\sin\theta$: the three logarithmic terms cancel by
$\sin\theta=2\sin(\theta/2)\cos(\theta/2)$.
Both sides tend to zero at the origin, proving the identity on
$(-1,1)$.  Finally $\operatorname{arctanh}(\sin x)$ has only an
integrable logarithmic singularity at $x=\pi/2$; dominated convergence
and oddness give the endpoint cases.
\end{proof}

The failure at a specific admissible negative parameter is exact.  At
$\lambda=-1/2$, $\theta=-\pi/6$, the printed principal-logarithm
right side differs from the integral by
\[
 i\pi\theta=-\frac{i\pi^2}{6}.
\]
The real integral is approximately
$-0.5317299165586139653837501727942413753$.
Thus the correction is a branch issue, not a numerical uncertainty.
The smallest repair is to insert the absolute value in the logarithm
and specify the continuous value at zero.  Restricting the printed
identity to $0<\lambda<1$ would also make its original logarithm valid.

\subsection{Remove a remaining unsupported non-elementarity label}

The same chapter, at line 194 of the pinned source, calls
\[
 \int_0^{\pi/2}\arctan^2(\sin x)\,dx
       =\pi\chi_2(3-2\sqrt2)
\]
the ``simplest non-elementary diagonal.''  The exact evaluation is
retained.  The manuscript gives neither a specified comparison algebra
nor a proof that this constant lies outside one.  Its later discussion
correctly distinguishes exact formulas from unproved minimality and
non-elementarity.  The proposed patch replaces the unsupported adjective
by an explicit statement that no elementary reduction or nonreduction
theorem for this value is established there.  This is a correction of
the claim's status; it is not a proof of a new reduction.

\subsection{Coefficient fields depend on the coordinates}

The paragraph following \texttt{integral:res:parity} in
\path{chapters/07-integration.tex} says that the reduction coefficients
live in $\mathbb Q(\sqrt q)$.  That description fits the immediately
displayed real quadratic character modulo five.  It should not be read
as a field assertion for arbitrary rational grids or arbitrary
characters.  Chapter~2 already correctly uses the cyclotomic DFT field
$\mathbb Q(\zeta_q)$; passing to multiplicative characters can also
introduce character values and Gauss sums.

For example, the primitive character modulo five specified by
$\chi(2)=i$ takes the values $1,i,-i,-1$ on $1,2,3,4$.
The character-coordinate coefficients therefore include $i$, which is
not in the real field $\mathbb Q(\sqrt5)$.  Even a real sine-coordinate
presentation need not stay in that quadratic field:
\[
 \sin^2\frac{2\pi}{5}=\frac{5+\sqrt5}{8},\qquad
 N_{\mathbb Q(\sqrt5)/\mathbb Q}
      \left(\sin^2\frac{2\pi}{5}\right)=\frac5{16}.
\]
If $\sin(2\pi/5)$ belonged to $\mathbb Q(\sqrt5)$, the last
quantity would be the square of a rational norm, which it is not.
The displayed quadratic-character identity itself requires no repair.
The appropriate clarification is to state the coefficient field for
the chosen additive, character, or real-paired coordinates.

\subsection{Prior in-repository continuations must be credited}

The pinned repository already contains substantial periodic Stieltjes
calculus outside the manuscript chapters, under
\path{docs/reports/stieltjes-correlation-calculus/continuations/}
relative to \path{Analysis/Polylogarithms/} \cite{RepoConvolutionAlgebra,RepoConvolutionCalculus,RepoShiftedJets}.
In particular, \path{convolution-algebra/article.tex} proves the
one-generator Bell algebra in \texttt{thm:Bell}, the Gamma-quotient
multiplication law in \texttt{thm:product}, the polygamma contact
correction, and the reflected spectral kernel in
\texttt{thm:reflected}.  The file
\path{convolution-calculus/Stieltjes_Convolution_Calculus.tex}
already records the symmetric first-Stieltjes order-zero
polylogarithm formula in \texttt{eq:Li0identity}.
The file \path{shifted-hurwitz-jets/sections.tex} gives the complete
normalized primitive tower in \texttt{shj:Uformula} and proves its
endpoint regularity.

The untwisted periodic results identified above therefore consolidate and
independently verify those untwisted foundations. The nonintegral-twist
section separately answers the explicit open target in the archived
convolution-algebra report.  The appearance of a
question in an earlier intake report does not establish its current open
status when a later in-repository continuation has already answered it.
The shifted harmonic Laurent and nonintegral-twist results are assessed
separately, with their own source and literature qualifications, rather
than deriving a novelty claim from the untwisted periodic construction.

An explicitly open target in the same
\path{convolution-algebra/article.tex} is the subsection
``Twists, characters, and alternate endpoint normalizations.''
It asks for the nonintegral frequency-shift analogue and an all-index
half-twist table with its point-supported terms.
Section~\ref{tw:sec:main} answers this specified target: the convolution
unit is $\delta_0$, the scalar Gamma coefficients remain ordinary zeta
values, the basis consists of Lerch jets, and the covariant contact
terms are given explicitly by Theorem~\ref{tw:thm:contacts}.
The nonsingular alternating convolution and the modified Stieltjes
functions themselves have the antecedents cited there.  The resolved
repository question concerns their singular extension and normalization.

\subsection{What this continuation does not prove}

At the pinned snapshot, the mixed sums $S_2$ and $S_4$ already have
exact proofs.  The weight-seven relation
\texttt{cycloquot:conj:S6} and the newer weight-nine relation
\texttt{s8new:conj:S8} remain conjectural.  A different frozen
weight-nine vector was rigorously rejected in
\texttt{research:prop:S8-rejected}; its rejection must not be confused
with the status of the newer candidate.  No result of the present
continuation promotes either surviving numerical candidate to a theorem.

The distributional identities proved above fix their extensions through
meromorphic continuation, or equivalently the stated cutoff finite parts.
Agreement with an ordinary function away from the endpoint and zero mean
alone do not characterize every higher derivative extension: derivatives
of the endpoint delta distribution have zero mean.  The explicit endpoint
terms in those theorems are part of the statements and must be preserved
when the formulas are integrated or differentiated.

% END input


% BEGIN sections/07_verification_questions.tex
\section{Verification, integration, and further research}
\label{future:sec}

\subsection{What has been proved and what has been checked}

The proofs in this article establish the identities. Their main
analytic steps are normal convergence after pairing Fourier
coefficients with a smooth test function, explicit Taylor subtraction
at a singular endpoint, joint meromorphic continuation before taking
parameter coefficients, and finite Gamma/Bernoulli coefficient
algebra. None of the theorems is inferred from a numerical fit.

The most delicate points were checked independently of their initial
derivations. These include the constant Fourier mode, the signs of
the covariant Dirac derivatives, the minus sign on the wrapped
half-twist interval, the harmonic Mellin remainder at a Gamma zero,
the depth-mixing reflection shift \(a+u\), and the terminal
\(\mathcal A_{-1}=1/\Gamma(1+u)\) term in rational multiplication.
The full mathematical review is supplied in
\path{verification/independent_review.md}. This is an independent
derivation within the present research workflow, not a claim of
external peer review or proof-assistant formalization.

The exact and numerical replays have different roles.
Exact symbolic checks verify finite polynomial consequences of the
stated generating identities. Numerical checks compare independent
analytic representations, with endpoint cancellations removed before
evaluation. They are floating-point diagnostics, not rigorous
interval enclosures.

\begin{center}
\begin{tabularx}{\textwidth}{@{}p{0.29\textwidth}X@{}}
\toprule
Verification component & Recorded coverage\\
\midrule
Untwisted coefficient algebra &
Six explicit Stieltjes products; reflected phase coefficients through
degree eight for both signs; 128 polygamma product cases and
30 repeated-primitive recurrences.\\[3pt]
Harmonic exact algebra &
167 checks: local kernels, Hurwitz--Bernoulli reductions, residue
differentiation, reflection, primitives, multiplication for moduli
two and three, and all displayed specializations.\\[3pt]
Periodic quadrature &
15 checks at 50 decimal digits, including finite-part test-function
pairings and ordinary digamma/log-Gamma convolutions.\\[3pt]
Harmonic Mellin continuation &
Three regular-value checks and four Laurent cases at each of
\(\varepsilon=10^{-4},10^{-5},10^{-6}\), using an independent scalar
continued Mellin evaluator at 110 decimal digits.\\[3pt]
Half-twist and source repair &
Nine checks at 60 decimal digits: the convergent half-twist integral
at three shifts, its quarter-Gamma evaluation, four branch-corrected
arctanh integrals, and the exact spurious imaginary term.\\[3pt]
Twisted spectral normalization &
The separate twisted replay tests rational Lerch conversion,
the retained zero mode, and the covariant contact coefficients.
Its complete parameters and residuals are recorded with the script.\\
\bottomrule
\end{tabularx}
\end{center}

For the harmonic continuation, the scalar integrand is
\[
 \frac{e^{-(a+1)t}}{1-e^{-t}}\,
       \frac{[-\log(1-e^{-t})]^r}{r!}.
\]
The independent evaluator integrates its convergent local
logarithmic series on \((0,1/2)\), after analytic continuation of the
individual Mellin monomials, and performs direct quadrature on
\((1/2,\infty)\). It does not evaluate through the
\(\mathcal A_m\) coefficient theorem. Degrees 80 and 100 agreed to
better than \(1.7\times10^{-92}\) in the recorded tests. After
subtracting the asserted principal and constant terms, the remainder
was \(O(\varepsilon)\) in every tested case. That agreement diagnoses
the coefficient extraction; it does not certify a truncation bound
for arbitrary inputs.

The extra half-twist quadrature uses only digamma and polygamma
functions on the integral side, and Hurwitz Stieltjes constants or
\(\log\Gamma(1/4)\) on the comparison side. A Taylor quotient removes
the cancellation at \(x=0\). Its largest recorded scaled residual,
including the branch-audit checks, is below \(4\times10^{-61}\).
Here a scaled residual means
\[
 \frac{|\hbox{computed left side}-\hbox{computed right side}|}
      {1+|\hbox{computed right side}|}.
\]
A tiny residual is useful implementation evidence, not an equality proof.

\subsection{Suggested repository integration}

The package is designed for archival intake first and mathematical
integration second. The supplied source correction patch applies
against the pinned chapter bytes. The main article has modular
sources and an equivalent standalone source, plus complete replay
scripts and recorded results.

The integration should preserve the following distinctions:
\begin{enumerate}[leftmargin=1.6em]
\item Cite the existing untwisted continuations instead of replacing
      their provenance with the present rederivation.
\item Add the nonintegral-twist section as a response to the archived
      half-twist question, retaining its cutoff definition, covariant
      derivative, unit \(\delta_0\), and reflection modulation.
\item Add the shifted harmonic Laurent theorem together with its
      precise outer-shift convention. Moving the inner harmonic
      denominators as well defines a different family.
\item Keep each finite-part identity attached to its domain and
      normalization. An isolated ordinary formula does not determine
      its endpoint distribution.
\item Apply the two narrowly scoped source corrections separately
      from accepting the new research statements.
\end{enumerate}
The machine-readable claim ledger records proved additions,
consolidated antecedents, corrections, and unresolved targets.
The proposed patch does not change the status of \(S_6\) or \(S_8\).

\subsection{Further research questions}

The following questions concern exact identities and reductions.
They are proposals for further work, not claims supported merely by
high-precision numerical evidence.

\begin{question}[Products with genuinely different twists]
\label{future:q:unequaltwists}
For real nonintegral \(\theta,\eta\), with
\(\theta-\eta\notin\mathbb Z\), determine a useful finite or
convergently evaluable normal form for
\[
 \widehat{\mathcal K_{\alpha,\beta}^{\theta,\eta}}(n)
 =(2\pi i(n+\theta))^\alpha
  (2\pi i(n+\eta))^\beta .
\]
Which parameter configurations permit reduction to finitely many
single-twist Lerch jets, and when is a genuinely two-shift
special function necessary?
\end{question}

The equal-twist case is now solved by the Gamma quotient.
The half-twist involution reverses frequencies within one twist;
it does not solve the unequal-twist problem. A successful extension
should retain its Laurent endpoint terms and rational-character
coordinates, rather than only its nonsingular Fourier expression.

\begin{question}[Regular Laurent coefficients beyond the finite part]
\label{future:q:regularcoefficients}
For \(j\geq1\), determine exact representations and reduction rules
for
\[
 [\varepsilon^j]E_r(-m+\varepsilon;a).
\]
At rational \(a\), when can these coefficients be expressed using
ordinary Hurwitz or Dirichlet-\(L\) derivatives, and when do new
height-one or colored multiple-zeta derivatives remain?
\end{question}

The Gamma zero eliminates the analytic Mellin remainder through
the constant coefficient, but not beyond it. Thus the polynomial
theorem proved here gives the full singular and finite data without
settling the next regular coefficient. The independent scalar
Mellin representation supplies a rigorous starting point for
such a study.

\begin{question}[Lifting the reflection and multiplication laws]
\label{future:q:lift}
Find exact formulas for the analytic defects obtained when
\eqref{harm:eq:FPreflection} and \eqref{harm:eq:multiplication}
are lifted from Laurent coefficient polynomials to the full
meromorphic family \(E_r(s;a)\).
Can the defects be organized as convergent iterated integrals or
colored polylogarithms with a natural depth filtration?
\end{question}

The finite polynomial identities are exact, not asymptotic.
Nevertheless their proof uses only the local Mellin kernel. A
global lift must also control the rest of the integral. Keeping
that local/global separation explicit should prevent a false
unmodified Hurwitz multiplication rule for the harmonic family.

\begin{question}[Rational-twist arithmetic at singular orders]
\label{future:q:characters}
Combine the all-index twisted table with the manuscript's
conductor-sensitive character conversion. Determine economical
primitive-character formulas for the resulting rational-shift
convolutions, including every conductor logarithm and endpoint
term. Which combinations descend to a proper sub-conductor?
\end{question}

A finite coordinate algorithm is already available; the remaining
problem is its exact arithmetic simplification. Reduction to a
chosen spanning set is a meaningful target even when independence
of the evaluated constants is unknown.

\begin{question}[Translated cubic and higher log-Gamma products]
\label{future:q:moments}
With \(g(x)=\log\Gamma(x)-\tfrac12\log(2\pi)\), obtain exact reductions
for the two-parameter ordinary product
\[
 C_3(a,b)=\int_0^1g(x)g(\{x+a\})g(\{x+b\})\,dx,
 \qquad 0<a<b<1,
\]
and for higher moments beyond the known unshifted cubic formula.
At rational \(a,b\), which colored multiple-zeta derivatives remain
after exact character and symmetry reductions?
\end{question}

The unshifted cubic moment is already evaluated in the pinned chapters
\path{07-loggamma-moments.tex} and \path{07-tornheim-evaluation.tex},
and must not be proposed as an unsolved problem
\cite{RepoManuscript}. The target here retains two distinct translation
parameters. These are pointwise products, distinct from convolution
powers: Fourier multiplication introduces several independent indices.
A concrete starting family is the meromorphic continuation of
\[
 \mathcal T(s,t,u;z,w)=
 \sum_{m,n\geq1}\frac{z^mw^n}{m^s n^t(m+n)^u},
\]
and its order derivatives at the integer triples arising from
Kummer's Fourier coefficients, with phases fixed by \(a,b\).
The research target is an explicit
reduction theorem, not a numerical assertion that a selected
constant basis is complete or independent.

\begin{question}[Collision regularizations across different products]
\label{future:q:collisions}
When translated singularities collide, compare the finite part
obtained from the meromorphic convolution family with finite parts
obtained by coincident-point cutoff procedures. Determine the
exact local counterterm relating these prescriptions, and which
associativity or differentiation identities each preserves.
\end{question}

The distribution convolution itself is defined on the compact
circle. Evaluating it at its singular point is a different
operation and needs a separate prescription. The zero-mode
degeneration theorem controls a twist resonance, not every
possible collision regularization of pointwise products.

\begin{question}[The remaining Gaussian candidates]
\label{future:q:targets}
Can the level-four iterated-integral and conductor-jet reductions
in the manuscript be completed into an exact certificate for
\texttt{cycloquot:conj:S6} or \texttt{s8new:conj:S8}?
If a conjectured coefficient vector is false, can it be rejected
with a certified error enclosure and replaced by a precisely
stated new target?
\end{question}

The source already distinguishes these candidates from the proved
\(S_2,S_4\) identities and the rejected earlier \(S_8\) vector.
The present twist and harmonic identities should be used as
additional transformations, without suggesting that their
existence proves either outstanding relation.

\paragraph{A practical formalization programme.}
A useful accompanying project is to formalize the finite coefficient
algebra separately from the analytic continuation lemmas. The exact
Bell, Bernoulli, reflection, multiplication, and residue recurrences
give compact symbolic certificates at any requested depth. The
analytic layer should then discharge the local-uniform convergence,
endpoint subtraction, and uniqueness hypotheses once, instead of
replacing them by growing collections of numerical examples.


% END input


% BEGIN references.tex
\begin{thebibliography}{99}
\addcontentsline{toc}{section}{References}
\small\raggedright

\bibitem{RepoManuscript}
ProveIt project,
\emph{Polylogarithm manuscript},
\path{Analysis/Polylogarithms/docs/manuscript/},
snapshot \pinned\ (10 October 2026).
\href{https://github.com/VladimirReshetnikov/ProveIt/tree/e0d9463bdee9685dfb1dddb819059cc738540c57/Analysis/Polylogarithms/docs/manuscript}{Pinned source directory}.
Exact inspected file hashes are supplied in \path{integration/source_manifest.json}.

\bibitem{RepoIncoming}
ProveIt project,
\emph{Polylogarithm--Stieltjes Continuation},
pending archive dated 10 October 2026,
\path{docs/incoming/ProveIt_Polylogarithm_Stieltjes_Continuation_2026-10-10.zip},
at the same snapshot.
Git blob \texttt{2b19b137342f68561fe7bca55da16caa2291baf1}.
\href{https://github.com/VladimirReshetnikov/ProveIt/tree/e0d9463bdee9685dfb1dddb819059cc738540c57/docs/incoming}{Pinned incoming directory}.

\bibitem{RepoConvolutionAlgebra}
ProveIt project,
\emph{A Convolution Algebra of Stieltjes Functions},
archived research continuation, 10 October 2026.
\path{stieltjes-correlation-calculus/continuations/convolution-algebra/article.tex},
under \path{Analysis/Polylogarithms/docs/reports/}.
\href{https://github.com/VladimirReshetnikov/ProveIt/blob/e0d9463bdee9685dfb1dddb819059cc738540c57/Analysis/Polylogarithms/docs/reports/stieltjes-correlation-calculus/continuations/convolution-algebra/article.tex}{Pinned source}.
The final section explicitly proposes the all-index half-twist
multiplication table and its point-supported counterterms.

\bibitem{RepoConvolutionCalculus}
ProveIt project,
\emph{A Convolution Calculus for Stieltjes Functions},
archived research continuation, 10 October 2026.
\path{stieltjes-correlation-calculus/continuations/convolution-calculus/Stieltjes_Convolution_Calculus.tex},
under \path{Analysis/Polylogarithms/docs/reports/}.
\href{https://github.com/VladimirReshetnikov/ProveIt/blob/e0d9463bdee9685dfb1dddb819059cc738540c57/Analysis/Polylogarithms/docs/reports/stieltjes-correlation-calculus/continuations/convolution-calculus/Stieltjes_Convolution_Calculus.tex}{Pinned source}.

\bibitem{RepoShiftedJets}
ProveIt project,
\emph{Shifted Hurwitz-Zeta Correlations:
Exact Stieltjes Closure, Polygamma Contact Terms, and Log-Gamma Antiderivatives},
archived continuation, 10 October 2026.
\path{stieltjes-correlation-calculus/continuations/shifted-hurwitz-jets/sections.tex},
under \path{Analysis/Polylogarithms/docs/reports/}.
\href{https://github.com/VladimirReshetnikov/ProveIt/blob/e0d9463bdee9685dfb1dddb819059cc738540c57/Analysis/Polylogarithms/docs/reports/stieltjes-correlation-calculus/continuations/shifted-hurwitz-jets/sections.tex}{Pinned source}.

\bibitem{RepoAlternating}
ProveIt project,
\emph{Reflection, certified evaluation, and a depth--exponent
transition for alternating harmonic polylogarithms},
archived continuation, 7 October 2026.
\path{Analysis/Polylogarithms/docs/reports/alternating-harmonic-polylogarithms/article.tex}.
\href{https://github.com/VladimirReshetnikov/ProveIt/blob/e0d9463bdee9685dfb1dddb819059cc738540c57/Analysis/Polylogarithms/docs/reports/alternating-harmonic-polylogarithms/article.tex}{Pinned source}.

\bibitem{RepoHarmonicGeometry}
ProveIt project,
\emph{Harmonic order geometry}, archived continuation.
\path{lerch-global-phase/continuations/harmonic-order-geometry/sections/harmonic_order.tex},
under \path{Analysis/Polylogarithms/docs/reports/}.
\href{https://github.com/VladimirReshetnikov/ProveIt/blob/e0d9463bdee9685dfb1dddb819059cc738540c57/Analysis/Polylogarithms/docs/reports/lerch-global-phase/continuations/harmonic-order-geometry/sections/harmonic_order.tex}{Pinned source}.
Used for the pre-existing scalar Mellin-subtraction mechanism.

\bibitem{DLMFHurwitz}
NIST Digital Library of Mathematical Functions,
\S25.11, Hurwitz Zeta Function.
\url{https://dlmf.nist.gov/25.11}.

\bibitem{DLMFBeta}
NIST Digital Library of Mathematical Functions,
\S5.12, Beta Function.
\url{https://dlmf.nist.gov/5.12}.

\bibitem{DLMFPolylog}
NIST Digital Library of Mathematical Functions,
\S25.12, Polylogarithms.
\url{https://dlmf.nist.gov/25.12}.

\bibitem{DLMFLerch}
NIST Digital Library of Mathematical Functions,
\S25.14, Lerch's Transcendent.
\url{https://dlmf.nist.gov/25.14}.

\bibitem{DLMFDirichlet}
NIST Digital Library of Mathematical Functions,
\S25.15, Dirichlet \(L\)-functions.
\url{https://dlmf.nist.gov/25.15}.

\bibitem{per:Sun}
Zhi-Hong Sun,
\emph{On the properties of invariant functions},
Bulletin des Sciences Math\'ematiques \textbf{189} (2023), 103347.
\url{https://arxiv.org/abs/2209.14625}.
The normalized Hurwitz convolution is in Remark 3.2, p.~15,
of the accessible arXiv v1.

\bibitem{per:GaySebbar}
R. Gay and A. Sebbar,
\emph{Pseudo-differential operators on the circle, Bernoulli polynomials},
Quantum Studies: Mathematics and Foundations \textbf{11} (2024), 1--25.
\url{https://doi.org/10.1007/s40509-024-00316-9}.

\bibitem{Kim2026}
Namhoon Kim,
\emph{Analytic Continuation of the Hurwitz Transform},
Mathematics \textbf{14}(2) (2026), article 271.
\url{https://doi.org/10.3390/math14020271}.

\bibitem{Young2023}
Paul Thomas Young,
\emph{Global series for height 1 multiple zeta functions},
European Journal of Mathematics \textbf{9} (2023), article 99.
\url{https://doi.org/10.1007/s40879-023-00695-0}.
The inspected publisher abstract explicitly states results on singular
Laurent parts and linear Laurent coefficients. A full-text theorem-by-theorem
comparison was not available.

\bibitem{Young2024}
Paul Thomas Young,
\emph{Series of Height One Multiple Zeta Functions},
Integers \textbf{24} (2024), A43.
\url{https://math.colgate.edu/~integers/y43/y43.pdf}.
DOI: \url{https://doi.org/10.5281/zenodo.11221606}.

\bibitem{Young2026}
Paul Thomas Young,
\emph{Parametric Euler Sums of Generalized Hyperharmonic Numbers},
Integers \textbf{26} (2026), A81.
\url{https://math.colgate.edu/~integers/aa81/aa81.pdf}.
DOI: \url{https://doi.org/10.5281/zenodo.20931235}.
Theorem 1 supplies the beta-derivative evaluation framework.

\bibitem{HarmonicStieltjes2023}
Levent Karg{\i}n, Ayhan Dil, Mehmet Cenkci, and M\"um\"un Can,
\emph{On the Stieltjes constants with respect to harmonic zeta functions},
Journal of Mathematical Analysis and Applications \textbf{525}
(2023), 127302.
\url{https://arxiv.org/abs/2304.03517}.
DOI: \url{https://doi.org/10.1016/j.jmaa.2023.127302}.

\bibitem{HuKim2022}
Su Hu and Min-Soo Kim,
\emph{On the Stieltjes constants and gamma functions with respect to
alternating Hurwitz zeta functions},
Journal of Mathematical Analysis and Applications \textbf{509}(1)
(2022), 125930.
\url{https://arxiv.org/abs/2106.14674}.
DOI: \url{https://doi.org/10.1016/j.jmaa.2021.125930}.

\bibitem{ZhuHuKim2026}
Haiqing Zhu, Su Hu, and Min-Soo Kim,
\emph{On the properties of alternating invariant functions},
arXiv:2505.10985v3 (9 September 2025).
\url{https://arxiv.org/pdf/2505.10985}.
Theorem 1.13 is the nonsingular alternating Hurwitz convolution.
Journal version: Bulletin des Sciences Math\'ematiques \textbf{206}
(2026), 103724,
\url{https://doi.org/10.1016/j.bulsci.2025.103724}.

\end{thebibliography}


% END input

\end{document}
